\chapter{Errors, Validation \& Problem Details (RFC 9457)}
\label{ch:problem_details}

\section*{Executive Summary}

Every API has two interfaces: the success contract and the failure contract.
Teams routinely invest weeks polishing the former and minutes on the latter,
and then discover---during a 3~a.m.\ incident---that clients cannot tell a
validation mistake from a melted database, cannot retry safely, and cannot
page the right owner because every error looks like \texttt{500 Internal
Server Error} with a stack trace taped to it. This chapter fixes that.

We treat errors as a first-class design surface. You will learn a rigorous
taxonomy that separates transport, protocol, and application failures, and a
repeatable decision procedure that classifies each failure as client-fixable
(4xx), server-owned (5xx), or dependency-caused, and as retriable or
non-retriable. You will master RFC~9457 \emph{Problem Details for HTTP APIs}
\cite{rfc9457}, the standard that replaces ad-hoc error payloads with a
single, extensible, machine-readable shape carried in
\texttt{application/problem+json}. You will learn status-code precision---why
\texttt{422} means \emph{syntactically valid but semantically invalid}, why
\texttt{404} and \texttt{410} tell different stories, why
\texttt{200}-with-an-error-body is an anti-pattern that corrupts caches,
retries, and monitoring.

On the engineering side, you will build a production-grade error pipeline in
FastAPI: strict Pydantic~v2 validation at the boundary
\cite{pydanticdocs}, a \texttt{ProblemException} hierarchy mapped to RFC~9457
responses by a single handler, PII-redaction middleware, correlation IDs
end-to-end, a machine-readable error-code registry, and a contract test suite
that schema-validates \emph{every} error path with JSON Schema
\cite{jsonschemadocs}. Everything runs offline, deterministically, in
process. By the end, ``what did the API say when it failed?'' will have one
answer everywhere in your system: a well-formed problem document.

\begin{keyconceptbox}
\textbf{Errors are part of the contract.} Version them, document them, test
them, and alert on them exactly as you do for success responses. A consumer
should be able to handle every failure mode of your API without ever reading
a log line or calling your support desk.
\end{keyconceptbox}

\section{Learning Objectives \& Prerequisites}

\subsection*{Learning Objectives}

After completing this chapter, you will be able to:

\begin{enumerate}
  \item Classify any failure into the transport/protocol/application taxonomy
        and justify a status code, retriability verdict, and ownership
        assignment for it.
  \item Author RFC~9457 problem documents, including extension members, and
        explain the semantics of \texttt{type}, \texttt{title},
        \texttt{status}, \texttt{detail}, and \texttt{instance}.
  \item Choose precisely between \texttt{400} and \texttt{422}, \texttt{404}
        and \texttt{410}, \texttt{401}, \texttt{403} and \texttt{407}, and
        among \texttt{500}, \texttt{502}, \texttt{503}, and \texttt{504}.
  \item Engineer validation with Pydantic~v2: strict mode,
        \texttt{field\_validator} versus \texttt{model\_validator},
        \texttt{Annotated} custom types, and coercion pitfalls.
  \item Implement a single error-handling pipeline---domain exceptions
        $\rightarrow$ problem mapper $\rightarrow$ redaction $\rightarrow$
        response---with correlation IDs and server-side structured logging.
  \item Schema-validate every error response in contract tests and generate
        a machine-readable error-code registry and its public documentation.
  \item Operate errors in production: error-rate SLIs per code and endpoint,
        alert thresholds, PII redaction, and error-budget interaction
        (Chapter~18).
\end{enumerate}

\subsection*{Prerequisites}

This chapter assumes you have completed Chapters~0--6. In particular, you
should be comfortable with:

\begin{itemize}
  \item HTTP semantics: methods, status-code classes, header negotiation
        (Chapter~1, \cite{rfc9110}).
  \item Consuming APIs from Python with \texttt{httpx} (Chapter~2).
  \item REST resource modeling and representation design (Chapter~3).
  \item Serialization contracts, JSON Schema, and OpenAPI~3.1 documents
        (Chapter~4, \cite{openapi31}).
  \item Client-side security hygiene: where secrets may and may not appear
        (Chapter~5).
  \item Building FastAPI services: routers, dependency injection, Pydantic
        models, and \texttt{TestClient} (Chapter~6, \cite{fastapidocs}).
\end{itemize}

All code targets Python~3.11+, runs fully offline, and uses only the
dependencies in the Python Dependencies Appendix. Commands are given for
Windows PowerShell.

\section{Deep Technical Mechanics \& Architectural Concepts}

\subsection{A Taxonomy of Failure}

Before choosing a status code or a payload shape, you must know
\emph{what kind of thing went wrong}. We use a three-layer taxonomy.

\begin{definition}[Transport Error]
A failure below HTTP semantics: the bytes never formed a complete, routable
request or response. Examples: TCP connection refused, TLS handshake failure,
DNS resolution failure, connection reset mid-body. Transport errors have
\emph{no HTTP status code at all}---there is nothing to attach one to. From
the client's perspective they surface as exceptions (e.g.\
\texttt{httpx.ConnectError}), and they are almost always retriable with
backoff because no application state was necessarily touched.
\end{definition}

\begin{definition}[Protocol Error]
A failure of HTTP itself: the message arrived but is malformed or violates
protocol expectations. Examples: syntactically invalid JSON in the request
body, a \texttt{Content-Type} the server does not support (\texttt{415}),
a missing mandatory header, an unacceptable \texttt{Accept} header
(\texttt{406}). The server understood nothing reliably enough to run business
logic, so the response must not imply that it did.
\end{definition}

\begin{definition}[Application (Domain) Error]
A failure inside valid protocol traffic: the request was well-formed HTTP
carrying well-formed JSON, but the domain rejected it. Examples: a SKU that
does not exist, a state transition the shipment lifecycle forbids,
insufficient funds. Application errors are where APIs earn or lose trust,
because they are the errors clients handle programmatically.
\end{definition}

The taxonomy matters because each layer has different handling machinery.
Transport errors are handled by client retry policies and circuit breakers
\cite{nygard2018releaseit}. Protocol errors are handled by the framework's
request parsing. Application errors are handled by \emph{your} domain code,
and they deserve the most design effort.

\subsubsection{Ownership: Who Can Fix It?}

The 4xx/5xx split in RFC~9110~\cite{rfc9110} is fundamentally about
\emph{ownership of the fix}:

\begin{itemize}
  \item \textbf{Client-fixable (4xx).} The client sent something the server
        can never accept as-is. Retrying the identical request is useless;
        something about the request must change. Examples: \texttt{400},
        \texttt{404}, \texttt{409}, \texttt{422}.
  \item \textbf{Server-owned (5xx).} The request was fine; the server failed
        to fulfill it. Retrying the identical request may succeed later.
        Examples: \texttt{500}, \texttt{503}.
  \item \textbf{Dependency-caused (also 5xx, but distinguishable).} The
        server failed \emph{because a downstream dependency failed}.
        \texttt{502} (bad response from upstream), \texttt{504} (upstream
        timeout), and \texttt{503} (deliberate unavailability) communicate
        this. Distinguishing dependency failures from internal bugs is
        operationally crucial: the page goes to a different team and the
        retry semantics differ.
\end{itemize}

\begin{warningbox}
\textbf{Never blame the client for your outage.} Returning \texttt{400} when
a database is down teaches clients to stop retrying, corrupts their error
metrics, and hides your incident behind their dashboards. If the request was
valid and you could not serve it, the answer is in the 5xx family---always.
\end{warningbox}

\subsubsection{A Decision Procedure for Retriability}

Clients need one bit of information above all: \emph{may I send this again?}
Encode the decision explicitly, for both client and server authors:

\begin{enumerate}
  \item \textbf{Did the request reach application logic?} If not (transport
        error, \texttt{408}, \texttt{429}, \texttt{502}, \texttt{503},
        \texttt{504}), the request is retriable \emph{if it is idempotent or
        idempotency-protected} (Chapter~10).
  \item \textbf{Is the failure about the request content?} \texttt{400},
        \texttt{401}, \texttt{403}, \texttt{404}, \texttt{409}, \texttt{415},
        \texttt{422}: not retriable without modification. \texttt{409} is the
        nuance: retriable \emph{after} the client refreshes state and
        re-evaluates preconditions.
  \item \textbf{Is the failure an unhandled server fault (\texttt{500})?}
        Treat as \emph{not} retriable for mutations unless the operation is
        idempotent; retrying a non-idempotent POST after a \texttt{500} risks
        double-execution because you cannot know whether the server
        committed before failing.
  \item \textbf{Did the server say when?} \texttt{429} and \texttt{503} with
        \texttt{Retry-After} are retriable \emph{at or after} the indicated
        time, with jitter to avoid thundering herds.
\end{enumerate}

Publish this classification per error code (we build exactly such a registry
in Listing~7.6) so no client engineer ever has to guess.

\subsection{RFC 9457: Problem Details for HTTP APIs}

Before RFC~7807 (2016), every API invented its own error payload:
\texttt{\{"error": "oops"\}}, \texttt{\{"message": \ldots, "code": 1234\}},
bare strings, HTML pages from middleware you forgot existed. Client code
became a museum of per-API parsers. RFC~7807 standardized a shape; RFC~9457
(March 2024) \cite{rfc9457} \emph{obsoletes} RFC~7807, clarifying extension
member handling, the \texttt{application/problem+xml} deprecation of
XML-specific guidance, and alignment with modern HTTP (RFC~9110
\cite{rfc9110}). If a document or library says ``7807,'' read it as 9457
with a historical accent.

\begin{definition}[Problem Detail Object]
A JSON (or XML) object carrying \emph{machine-readable details of one HTTP
error response}, served with the media type
\texttt{application/problem+json} and an error status code. It is a
deliberately minimal kernel---five standard members---extended by the API
with additional members.
\end{definition}

\subsubsection{The Five Standard Members}

\begin{table}[h]
  \centering
  \small
  \begin{tabularx}{\textwidth}{l l X}
    \toprule
    \textbf{Member} & \textbf{Required} & \textbf{Semantics} \\
    \midrule
    \texttt{type} & yes (default \texttt{about:blank}) & A URI reference that
      identifies the problem \emph{type}. When dereferenced, it should lead
      to human documentation. It is the stable identifier clients match on. \\
    \texttt{title} & recommended & Short, human-readable summary of the
      problem \emph{type}. Stable per type; never per-occurrence. Localizable
      but should not change between occurrences of the same type. \\
    \texttt{status} & recommended & The HTTP status code, echoed into the body
      for convenience of intermediaries and logs. Must match the actual
      response status. \\
    \texttt{detail} & optional & Human-readable explanation of \emph{this
      occurrence}. May differ per occurrence; never write client logic
      against it. \\
    \texttt{instance} & optional & A URI reference identifying \emph{this
      occurrence}; often the request path or a per-error URN. \\
    \bottomrule
  \end{tabularx}
  \caption{Standard problem detail members per RFC 9457 \cite{rfc9457}.}
  \label{tab:problem-members}
\end{table}

The division of labor is the design's genius: \texttt{type}/\texttt{title}
describe the \emph{class} of problem (stable, documentable, matchable);
\texttt{detail}/\texttt{instance} describe the \emph{occurrence} (specific,
variable, not for matching). Most broken client integrations come from
confusing these two levels---matching on \texttt{detail} strings is the
classic offender.

\subsubsection{Extension Members and Documentation-as-Contract}

RFC~9457 explicitly permits \emph{extension members}: additional keys defined
by your API, such as \texttt{code}, \texttt{trace\_id}, or an
\texttt{errors} array of per-field failures. Two rules keep extensions
healthy:

\begin{enumerate}
  \item Clients \textbf{must ignore unknown members}. This makes extension
        members additive-evolvable: you can ship new members without a
        version bump, exactly like new fields in a success representation.
  \item Every extension member you emit \textbf{must be documented at the
        \texttt{type} URI}. The \texttt{type} URI is the contract anchor:
        \emph{documentation-as-contract}. When a client dereferences
        \texttt{https://docs.example.invalid/problems/validation-failed},
        they should find the member list, the status code, retriability, and
        remediation guidance---generated from the same registry your code
        uses (Section~\ref{sec:impl}, Listing~7.6).
\end{enumerate}

RFC~9457 also defines an IANA registry of problem \emph{types}, though in
practice most APIs mint their own type URIs under a domain they control. The
operational requirement is not registration; it is \emph{stability}: a
\texttt{type} URI, once shipped, must keep meaning the same problem class
forever.

\subsubsection{Why Must \texttt{type} Be a URI?}

Three reasons. First, \emph{global uniqueness without coordination}: URIs
minted under your domain cannot collide with another vendor's
``\texttt{validation\_error}.'' Second, \emph{dereferenceability}: a URI can
double as a documentation link, collapsing identifier and manual into one
artifact. Third, \emph{opacity}: because the spec says clients match on the
URI string itself (not on fetched content), you may serve
\texttt{about:blank} for generic errors and reserve rich type URIs for
documented classes. Note that \texttt{about:blank} with a \texttt{title}
matching the status's reason phrase is the spec-sanctioned ``I have nothing
more specific to say'' fallback.

\subsubsection{Security Considerations}

RFC~9457's security section is blunt, and production experience blunter:
\textbf{problem details are exfiltration channels if you let them be}.
Rules we enforce in this chapter's code:

\begin{itemize}
  \item \textbf{No internals.} No stack traces, framework versions, SQL
        fragments, hostnames, or file paths in \texttt{detail}. Each is
        reconnaissance data for an attacker probing your error paths.
  \item \textbf{No PII.} \texttt{detail} must never echo submitted personal
        data. ``Email jane.doe@example.invalid already registered'' both
        leaks an account-existence oracle and lands PII in every log
        aggregator and APM tool that captures response bodies.
  \item \textbf{No secret material.} Failed-payload logging must redact
        credentials, tokens, and keys \emph{before} writing (Chapter~5).
  \item \textbf{Defense in depth.} Even with disciplined handlers, add a
        redaction middleware as the last gate (Listing~7.3): humans make
        mistakes at 3~a.m.; middleware does not.
\end{itemize}

\subsection{Status-Code Precision for Error Cases}

Status codes are the coarsest, most universally understood part of your error
contract; caches, load balancers, retry middleware, and monitoring all key on
them before anyone reads a body. Imprecision here is amplified everywhere.
This section walks the pairs that production systems most often confuse; all
semantics are per RFC~9110~\S 15 \cite{rfc9110}.

\subsubsection{400 versus 422}

\texttt{400 Bad Request} means the server \emph{cannot or will not process
the request due to something perceived to be a client error}---with the
canonical reading that the \emph{framing or syntax} itself is wrong:
malformed JSON, a missing required header, a query parameter that fails to
parse as an integer. The request could not be interpreted as a valid
instance of your contract.

\texttt{422 Unprocessable Content} means the request \emph{was} syntactically
valid---the JSON parsed, the schema shape held---but \emph{semantic}
instructions could not be followed: an end date before a start date, a
quantity of zero, a SKU failing its checksum. The mnemonic:
\textbf{400 = I cannot read you; 422 = I read you and reject what you said.}
FastAPI emits \texttt{422} for Pydantic validation failures by default
\cite{fastapidocs}, which is the convention this chapter keeps. Using
\texttt{400} for semantic failures is not fatal, but it is lossy: clients
distinguish ``fix my serializer'' from ``fix my business values'' only if
you do.

\subsubsection{404 versus 410}

\texttt{404 Not Found} says \emph{nothing} about history: the resource is not
here, and the server does not say whether it ever was. \texttt{410 Gone} is a
deliberate tombstone: the resource existed and is permanently retired---a
deleted account, an expired export, a sunset endpoint version. \texttt{410}
is cacheable by default and tells crawlers and sync clients to stop asking
and purge local copies. Use it when you know; use \texttt{404} otherwise.
Be aware that \texttt{410} (like any distinct response to existence probes)
can be an information-disclosure oracle for sensitive resources; for those,
uniform \texttt{404} is the privacy-preserving choice.

\subsubsection{409 Conflict and Its Domains}

\texttt{409} means the request is fine but \emph{conflicts with current
resource state}. It appears in at least four distinct domains, and your
problem \texttt{type} should say which:

\begin{itemize}
  \item \textbf{State-machine conflict:} cancelling an already-shipped
        shipment (our Listing~7.4 example).
  \item \textbf{Uniqueness conflict:} creating a resource whose natural key
        is taken. (Some APIs use \texttt{422} here; \texttt{409} is the
        better fit because the payload is valid in isolation.)
  \item \textbf{Concurrent-modification conflict:} an \texttt{ETag} /
        \texttt{If-Match} precondition failed---though \texttt{412} is the
        precise code for failed preconditions (Chapter~10).
  \item \textbf{Reference conflict:} deleting a warehouse that still holds
        inventory.
\end{itemize}

In all domains, \texttt{409} is \emph{conditionally retriable}: the client
should re-read state, re-decide, and only then retry.

\subsubsection{401 versus 403 versus 407}

\texttt{401 Unauthorized} is badly named: it means \emph{unauthenticated}---
``I do not know who you are (or your credentials failed).'' It
\emph{must} carry a \texttt{WWW-Authenticate} challenge header; a
\texttt{401} without it is a spec violation. \texttt{403 Forbidden} means
\emph{authenticated but not authorized}: ``I know who you are, and no.''
Re-authenticating will not help. \texttt{407 Proxy Authentication Required}
is the same challenge as \texttt{401} but issued by an intermediary proxy,
with \texttt{Proxy-Authenticate}; API authors rarely emit it, but gateway
operators must not strip it. One more subtlety: for resources whose
\emph{existence} is sensitive, authorization failure should collapse to
\texttt{404}, not \texttt{403}, to avoid an existence oracle.

\subsubsection{429 and Retry-After}

\texttt{429 Too Many Requests} says the client itself is the load problem
(Chapter~14 develops the limiting algorithms). It is only useful when paired
with \texttt{Retry-After} (delay-seconds or an HTTP date) and, ideally,
rate-limit headers describing the policy window. A \texttt{429} without
\texttt{Retry-After} invites clients to invent their own backoff---usually
badly, usually synchronized into thundering herds. In problem details, carry
the retry budget as an extension member so programmatic clients do not parse
headers at all.

\subsubsection{500, 502, 503, 504}

The 5xx family differentiates \emph{where} the failure lives:

\begin{itemize}
  \item \texttt{500 Internal Server Error}: your code faulted. Generic,
        uncatchable-by-design, always a bug or an unmodeled failure.
  \item \texttt{502 Bad Gateway}: a gateway/upstream returned something
        invalid. The dependency answered---wrongly.
  \item \texttt{503 Service Unavailable}: deliberate unavailability---
        overload shedding, maintenance, a blown circuit breaker
        \cite{nygard2018releaseit}. Almost always retriable; pair with
        \texttt{Retry-After} when the duration is known.
  \item \texttt{504 Gateway Timeout}: the dependency never answered in time.
\end{itemize}

The operational value is paging precision: \texttt{500} pages your team;
\texttt{502}/\texttt{504} page the integration owner; \texttt{503} should
mostly page nobody---it is the system protecting itself. If your error
mapper collapses all four into \texttt{500}, every dependency outage becomes
your incident.

\subsubsection{The 200-with-Error-Payload Anti-Pattern}

Some APIs return \texttt{200 OK} with \texttt{\{"success": false,
"error": \ldots\}} in the body. Every layer of the HTTP stack is then
actively misled: caches may store the failure as a success; retry middleware
sees success and does nothing; CDN and load-balancer metrics report 100\%
availability during a total outage; monitoring must parse bodies to see the
fire. The only defensible carve-outs are protocols that deliberately tunnel
errors above transport semantics (GraphQL over HTTP is the notable one,
Chapter~12)---and even there, the community is moving toward distinguishing
transport success from application errors explicitly. For REST-style HTTP
APIs: the status code must reflect the outcome. Always.

\begin{figure}[h]
\centering
\begin{tikzpicture}[
  font=\small,
  q/.style={draw, rounded corners, align=center, fill=blue!5, text width=3.4cm, inner sep=4pt},
  a/.style={draw, rectangle, align=center, fill=green!10, text width=2.5cm, inner sep=3pt},
  lbl/.style={note, midway},
  every edge/.style={draw, ->}
]
\node[q] (start) {Failure reached application logic?};
\node[q, below left=1.1cm and 0.4cm of start] (client) {Client-fixable?\\(request must change)};
\node[q, below right=1.1cm and 0.4cm of start] (server) {Dependency-caused?};
\node[q, below left=1.0cm and -0.2cm of client] (syntax) {Syntactically valid?};
\node[q, below right=1.0cm and -0.2cm of client] (domain) {State conflict or missing resource?};

\node[a, below left=0.9cm and -0.6cm of syntax] (c400) {400 malformed\\415 type};
\node[a, below right=0.9cm and -0.6cm of syntax] (c422) {422 semantic};
\node[a, below left=0.9cm and -0.6cm of domain] (c404) {404 / 410};
\node[a, below right=0.9cm and -0.6cm of domain] (c409) {409 conflict\\412 precondition};

\node[a, below left=0.9cm and -0.2cm of server] (s5x) {502 bad reply\\504 timeout};
\node[a, below right=0.9cm and -0.2cm of server] (s500) {500 own fault\\503 shed + Retry-After};

\path (start) edge node[lbl, above left=-2pt and -14pt] {yes} (client)
      (start) edge node[lbl, above right=-2pt and -14pt] {yes, 5xx} (server)
      (client) edge node[lbl, left] {content error} (syntax)
      (client) edge node[lbl, right] {resource/state error} (domain)
      (syntax) edge node[lbl, left] {no} (c400)
      (syntax) edge node[lbl, right] {yes} (c422)
      (domain) edge node[lbl, left] {missing/gone} (c404)
      (domain) edge node[lbl, right] {conflict} (c409)
      (server) edge node[lbl, left] {upstream} (s5x)
      (server) edge node[lbl, right] {internal} (s500);
\node[note, below=0.25cm of c404] {Auth overlays all: 401 challenge, 403 forbidden, 429 + Retry-After for throttling.};
\end{tikzpicture}
\caption{Status-code decision tree for API error responses. Auth failures
(401/403) and throttling (429) are orthogonal overlays: they can interrupt
any branch before application logic runs.}
\label{fig:status-tree}
\end{figure}

\subsection{Validation Engineering with Pydantic v2}

Validation is where error contracts meet user experience: a validation error
is the error a human being most often has to read, understand, and act on.
Pydantic~v2 \cite{pydanticdocs} gives us the machinery; discipline turns it
into a contract.

\subsubsection{Validation at Boundaries versus Domain Invariants}

\emph{Boundary validation} answers: is this payload a legal instance of the
input schema? Types, ranges, patterns, string lengths. It is mechanical,
declarative, and belongs in the model layer at the edge of the system.
\emph{Domain invariants} answer: is this payload legal \emph{given the
current state of the world}? ``Express shipments under 30~kg'' is an
invariant expressible in the model; ``this SKU is stocked at the destination
warehouse'' is an invariant that requires a database and therefore belongs
in the service layer, where its failure becomes a domain exception
(\texttt{409}/\texttt{422}), not a schema error. Keeping the two apart keeps
models reusable and keeps service errors semantically typed instead of
masquerading as typos.

\subsubsection{\texttt{field\_validator} versus \texttt{model\_validator}}

A \texttt{field\_validator} sees exactly one field (plus, if requested,
already-validated sibling values via \texttt{ValidationInfo}) and answers
single-field questions: is this SKU pattern valid, is this quantity within
review limits. A \texttt{model\_validator} with \texttt{mode="after"} sees
the whole, already-coerced model and answers cross-field questions: priority
versus total weight. Use field validators for local rules---they compose and
localize error pointers naturally---and reserve model validators for genuine
invariants. A model validator's error has no natural field location, which is
why our mapping emits the JSON Pointer \texttt{/} (the whole document) for
it.

\subsubsection{Coercion Strictness}

Pydantic's default \emph{lax} mode coerces: the string \texttt{"12"} becomes
the integer \texttt{12}; \texttt{"true"} becomes \texttt{True}. At a public
API boundary this is a hazard class: silent coercion converts client bugs
into corrupted records---the client \emph{thought} it sent a string, you
stored an integer, and their audit diff never matches. \emph{Strict mode}
(\texttt{model\_config = ConfigDict(strict=True)}) accepts only values of
the declared type, surfacing the type mismatch as a \texttt{422} the client
can fix. The nuance, demonstrated in Listing~7.2: JSON has no decimal
literal, so blanket strictness makes a \texttt{Decimal} field unusable. The
correct move is \emph{per-field} strictness opt-out
(\texttt{Field(strict=False, ...)}) hardened with a validator (finite,
quantized). Strictness is a dial per field, not a religion.

\subsubsection{Error Message UX}

A validation message is read by a tired integrator at 2~a.m. Good messages
name the rule and the expectation (``priority must be one of [express,
freight, standard]''), never echo the raw offending value (it may be a
credential they pasted by mistake), and never expose internal field or table
names that differ from the public contract. When you translate Pydantic's
error array into problem-details extensions, locate each failure with a JSON
Pointer (RFC~6901): \texttt{/lines/0/quantity} is unambiguous to both humans
and machines, and it survives localization of the human message.

\subsection{Error Envelope Architecture}

The cardinal rule: \textbf{one pipeline, one shape}. Every error that leaves
the service---validation failure, domain rejection, unhandled exception,
framework-generated error---passes through the same stages and emerges as
\texttt{application/problem+json}. Figure~\ref{fig:error-pipeline} shows the
pipeline our code implements.

\begin{figure}[h]
\centering
\begin{tikzpicture}[
  font=\small,
  stage/.style={draw, rounded corners, align=center, fill=blue!5, text width=2.9cm, inner sep=5pt},
  store/.style={draw, rectangle, align=center, fill=green!10, text width=3.1cm, inner sep=4pt},
  side/.style={note, align=center, text width=3.2cm},
  every edge/.style={draw, ->, thick}
]
\node[stage] (raise) {Domain code raises \texttt{ProblemException}};
\node[stage, right=0.7cm of raise] (map) {Exception handler maps to \texttt{ProblemDetail}};
\node[stage, right=0.7cm of map] (redact) {Redaction middleware scrubs PII / traces};
\node[stage, right=0.7cm of redact] (resp) {\texttt{application/problem+json} + status};

\node[store, below=0.9cm of map] (log) {Structured log: full context + stack trace keyed by \texttt{trace\_id}};

\path (raise) edge (map)
      (map) edge (redact)
      (redact) edge (resp)
      (map) edge (log);
\node[side, above=0.35cm of map] {correlation middleware stamps \texttt{trace\_id} first};
\node[side, below=0.35cm of redact] {client never sees internals; operators see everything};
\end{tikzpicture}
\caption{Single error-handling pipeline: exception $\rightarrow$ problem
mapper $\rightarrow$ redaction $\rightarrow$ response, with full-fidelity
server-side logging forked off at the mapper under the correlation ID.}
\label{fig:error-pipeline}
\end{figure}

\subsubsection{Machine-Readable Codes}

The RFC~9457 \texttt{type} URI is the primary machine identifier; we add a
\texttt{code} extension member (\texttt{NRW-SHIPMENT-NOT-FOUND}) because
operations teams alert, dashboard, and page on short codes far more easily
than on URIs, and because codes survive documentation-host moves. Codes must
be: unique, immutable once shipped, prefixed per system, and backed by a
registry that validates itself at import time (Listing~7.6). Never let a
developer invent a code inline in a handler---that is how
\texttt{ERR\_2} happens.

\subsubsection{Localization}

\texttt{title} and \texttt{detail} are human-facing and may be localized via
\texttt{Accept-Language}; \texttt{type} and \texttt{code} are machine-facing
and must \emph{never} be localized. The practical pattern: treat
\texttt{title} as a lookup key into a message catalogue per locale, keep
\texttt{code} ASCII, and keep interpolation values (limits, field names) as
extension members so localized strings are pure templates. Clients that
match on errors must match on \texttt{type}/\texttt{code} only---a localized
\texttt{title} is a rendering, not an identifier.

\subsection{Operational Concerns}

\subsubsection{Correlation IDs End-to-End}

Every request carries a trace identifier: accepted from an inbound
\texttt{X-Trace-Id} header if well-formed (bounded length, character
allowlist---never trust and reflect arbitrary header bytes), minted otherwise,
echoed in the response header, embedded in the problem body as
\texttt{trace\_id}, and attached to every log line the request produces.
This collapses the support loop from ``describe what you saw'' to ``give me
the trace\_id.'' Propagate the same ID to downstream calls so a single
customer report reconstructs the whole request tree (Chapter~18 generalizes
this to full distributed tracing).

\subsubsection{Error Observability and Alerting}

The error contract is your observability taxonomy: emit metrics as
\texttt{errors\_total\{endpoint, status, code\}} and you can answer, per
deploy, per endpoint, per error code: is this new, is it ours, is it a
dependency, is it client abuse. Suggested alert thresholds as a starting
point: 5xx rate $> 1\%$ of requests over 5 minutes pages; a single
\texttt{code} spiking $> 10\times$ its 7-day baseline pages; 4xx rate $>
25\%$ sustained opens a ticket (usually a broken client deploy, not your
bug). Error budget interaction is direct: 5xx responses burn the budget of
Chapter~18. Formally, over a window $W$:

\begin{equation}
\text{burn}(W) = \frac{\left|\{\,r \in W : \text{status}(r) \ge 500\,\}\right|}{\left|W\right|}
\qquad\text{budget}(W) = 1 - \text{SLO}
\end{equation}

and when $\text{burn}(W)$ approaches $\text{budget}(W)$, Chapter~18's policy
freezes feature deploys in favor of reliability work. A precise error
taxonomy is what makes that arithmetic actionable instead of theatrical.

\section{Production-Ready Code Implementation}
\label{sec:impl}

The reference implementation is a small FastAPI service for the fictional
\emph{Northwind Robotics} universe (all data synthetic). Six files form one
coherent system:

\begin{itemize}
  \item \texttt{problem\_details.py} --- the reusable RFC~9457 library:
        models, exception hierarchy, handlers, JSON Pointer mapping.
  \item \texttt{shipment\_models.py} --- strict Pydantic~v2 validation
        showcase with a coercion-pitfall demo.
  \item \texttt{redaction.py} --- correlation-ID and PII-redaction
        middleware.
  \item \texttt{app.py} --- the service wiring, with deliberately
        interesting failure endpoints.
  \item \texttt{test\_contract.py} --- contract tests schema-validating every
        error path.
  \item \texttt{error\_codes.py} --- the machine-readable error-code
        registry and its documentation generator.
\end{itemize}

Setup (PowerShell, offline after install):

\begin{minted}{text}
python -m venv .venv
.\.venv\Scripts\Activate.ps1
pip install -r requirements.txt   # see the Dependencies Appendix
pytest -q                         # 10 contract tests, all in-process
python shipment_models.py         # coercion pitfall demo
python error_codes.py             # regenerate the error catalogue
\end{minted}

\subsection{A Reusable Problem-Details Library}

Listing~7.1 is the heart of the chapter. \texttt{ProblemDetail} models the
RFC~9457 kernel plus our extension members (\texttt{code},
\texttt{trace\_id}, \texttt{errors[]}). \texttt{ProblemException} is the
single bridge between domain code and HTTP: domain code raises semantically
named subclasses; one handler renders them. The
\texttt{RequestValidationError} handler converts Pydantic's error array into
RFC~6901 JSON Pointers, and the catch-all \texttt{Exception} handler proves
the invariant that \emph{nothing} leaves the service unshaped---the client
gets a redacted \texttt{500} with a \texttt{trace\_id}; the log gets the
truth.

\begin{minted}{python}
"""problem_details.py -- Reusable RFC 9457 problem-details library.

Listing 7.1. Offline-first, stdlib + pydantic + fastapi only.
"""
from __future__ import annotations

import re
import uuid
from typing import Any, Optional

from fastapi import FastAPI, Request
from fastapi.exceptions import RequestValidationError
from fastapi.responses import JSONResponse
from pydantic import BaseModel, ConfigDict, Field

PROBLEM_MEDIA_TYPE = "application/problem+json"

# Synthetic, obviously fictional documentation namespace (RFC 9457 type URIs).
TYPE_BASE = "https://docs.example.invalid/problems"


class ErrorLocation(BaseModel):
    """One field-level validation failure, located by JSON Pointer (RFC 6901)."""

    pointer: str = Field(description="JSON Pointer to the offending value")
    code: str = Field(description="Machine-readable error code")
    message: str = Field(description="Human-readable, redacted message")


class ProblemDetail(BaseModel):
    """RFC 9457 problem detail object with Northwind Robotics extensions."""

    model_config = ConfigDict(populate_by_name=True)

    type: str = Field(default="about:blank")
    title: str
    status: int
    detail: Optional[str] = None
    instance: Optional[str] = None
    # Extension members.
    code: Optional[str] = None
    trace_id: Optional[str] = None
    errors: list[ErrorLocation] = Field(default_factory=list)


class ProblemException(Exception):
    """Base class: a domain exception that maps 1:1 to a ProblemDetail."""

    status: int = 500
    title: str = "Internal Server Error"
    type_slug: str = "internal"
    code: str = "NRW-INTERNAL"

    def __init__(
        self,
        detail: Optional[str] = None,
        *,
        instance: Optional[str] = None,
        errors: Optional[list[ErrorLocation]] = None,
    ) -> None:
        super().__init__(detail)
        self.detail = detail
        self.instance = instance
        self.errors = errors or []

    def to_problem(self, trace_id: str) -> ProblemDetail:
        return ProblemDetail(
            type=f"{TYPE_BASE}/{self.type_slug}",
            title=self.title,
            status=self.status,
            detail=self.detail,
            instance=self.instance,
            code=self.code,
            trace_id=trace_id,
            errors=self.errors,
        )


class ShipmentNotFound(ProblemException):
    status = 404
    title = "Shipment Not Found"
    type_slug = "shipment-not-found"
    code = "NRW-SHIPMENT-NOT-FOUND"


class ShipmentStateConflict(ProblemException):
    status = 409
    title = "Shipment State Conflict"
    type_slug = "shipment-state-conflict"
    code = "NRW-SHIPMENT-STATE-CONFLICT"


class WarehouseUnavailable(ProblemException):
    status = 503
    title = "Dependency Unavailable"
    type_slug = "warehouse-unavailable"
    code = "NRW-WAREHOUSE-UNAVAILABLE"


class RateLimitExceeded(ProblemException):
    status = 429
    title = "Rate Limit Exceeded"
    type_slug = "rate-limit-exceeded"
    code = "NRW-RATE-LIMIT-EXCEEDED"

    def __init__(self, retry_after: int, **kwargs: Any) -> None:
        super().__init__(
            detail=f"Retry after {retry_after} seconds.", **kwargs
        )
        self.retry_after = retry_after


def json_pointer(loc: list[Any]) -> str:
    """Convert a Pydantic error `loc` tuple to an RFC 6901 JSON Pointer."""
    parts = []
    for seg in loc:
        if seg == "body":  # FastAPI wraps body params under "body"
            continue
        parts.append(str(seg).replace("~", "~0").replace("/", "~1"))
    return "/" + "/".join(parts)


def problem_response(
    problem: ProblemDetail,
    headers: Optional[dict[str, str]] = None,
) -> JSONResponse:
    payload = problem.model_dump(exclude_none=True)
    return JSONResponse(
        status_code=problem.status,
        content=payload,
        media_type=PROBLEM_MEDIA_TYPE,
        headers=headers,
    )


def register_problem_handlers(app: FastAPI) -> None:
    """Single pipeline: domain exception -> problem mapper -> redaction -> response."""

    @app.exception_handler(ProblemException)
    async def handle_problem(
        request: Request, exc: ProblemException
    ) -> JSONResponse:
        trace_id = getattr(request.state, "trace_id", str(uuid.uuid4()))
        problem = exc.to_problem(trace_id)
        headers = None
        if isinstance(exc, RateLimitExceeded):
            headers = {"Retry-After": str(exc.retry_after)}
        return problem_response(problem, headers=headers)

    @app.exception_handler(RequestValidationError)
    async def handle_validation(
        request: Request, exc: RequestValidationError
    ) -> JSONResponse:
        trace_id = getattr(request.state, "trace_id", str(uuid.uuid4()))
        # 422: the payload was syntactically valid JSON but semantically invalid.
        locs = [
            ErrorLocation(
                pointer=json_pointer(list(err.get("loc", ()))),
                code=str(err.get("type", "invalid")),
                # Pydantic messages can echo raw input; drop ctx, keep it safe.
                message=str(err.get("msg", "Invalid value")),
            )
            for err in exc.errors()
        ]
        problem = ProblemDetail(
            type=f"{TYPE_BASE}/validation-failed",
            title="Validation Failed",
            status=422,
            detail="Request payload failed semantic validation.",
            instance=str(request.url.path),
            code="NRW-VALIDATION-FAILED",
            trace_id=trace_id,
            errors=locs,
        )
        return problem_response(problem)

    @app.exception_handler(Exception)
    async def handle_unexpected(
        request: Request, exc: Exception
    ) -> JSONResponse:
        trace_id = getattr(request.state, "trace_id", str(uuid.uuid4()))
        # Full context goes to the server log; the client gets a redacted body.
        print(f"trace_id={trace_id} unhandled={exc!r}")
        problem = ProblemDetail(
            type=f"{TYPE_BASE}/internal",
            title="Internal Server Error",
            status=500,
            detail="An unexpected error occurred. "
            "Reference the trace_id when contacting support.",
            instance=str(request.url.path),
            code="NRW-INTERNAL",
            trace_id=trace_id,
        )
        return problem_response(problem)


SECRET_PATTERNS: list[tuple[re.Pattern[str], str]] = [
    (re.compile(r"[\w.+-]+@[\w-]+\.[\w.]+"), "<email>"),
    (re.compile(r"\b\d{3}-\d{2}-\d{4}\b"), "<tax-id>"),
    (re.compile(r"(?i)(api[_-]?key|token|password)=[^&\s]+"), r"\1=<redacted>"),
]


def redact_text(text: str) -> str:
    """Scrub PII-like patterns and secrets from arbitrary text."""
    for pattern, replacement in SECRET_PATTERNS:
        text = pattern.sub(replacement, text)
    return text
\end{minted}

\begin{examplebox}
\textbf{A client handles this without string matching.} Given any error from
the shipments API, a Python client switches on the stable extension member:
\begin{minted}{python}
import httpx

resp = httpx.post("https://api.example.invalid/shipments", json=payload)
if resp.is_error:
    problem = resp.json() if resp.headers["content-type"].startswith(
        "application/problem+json") else {}
    match problem.get("code"):
        case "NRW-VALIDATION-FAILED":
            fix_fields(problem["errors"])          # JSON Pointers per field
        case "NRW-RATE-LIMIT-EXCEEDED":
            sleep(float(resp.headers["Retry-After"]))
        case _:
            escalate(problem.get("trace_id"))      # never match on "detail"
\end{minted}
Note what the client never does: parse \texttt{detail}, match English prose,
or guess retriability from the status alone.
\end{examplebox}

\subsection{Strict Validation Showcase}

Listing~7.2 models the Northwind Robotics \texttt{Shipment} resource. Study
the \texttt{Annotated} pattern: \texttt{Sku} and \texttt{CountryCode} are
named, reusable types whose constraints travel with the type itself---the
self-documenting alternative to scattering \texttt{Field(...)} magic at every
use site. Note also the deliberate \texttt{strict=False} opt-out on
\texttt{unit\_weight\_kg}: JSON has no decimal literal, so a strict
\texttt{Decimal} would reject \emph{every} legitimate JSON number; the
companion validator restores the safety (finiteness, quantization) that
strictness was buying.

\begin{minted}{python}
"""shipment_models.py -- Strict Pydantic v2 validation showcase.

Listing 7.2. Northwind Robotics "Shipment" resource. Demonstrates strict mode,
field_validator vs model_validator, and Annotated custom types.
"""
from __future__ import annotations

from decimal import Decimal
from typing import Annotated, Optional

from pydantic import (
    BaseModel,
    ConfigDict,
    Field,
    StringConstraints,
    field_validator,
    model_validator,
)

# --- Custom types via Annotated: reusable, composable, self-documenting. ----
Sku = Annotated[
    str,
    StringConstraints(pattern=r"^NRW-[A-Z]{2}-\d{4}$", min_length=11, max_length=11),
]
CountryCode = Annotated[str, StringConstraints(pattern=r"^[A-Z]{2}$")]


class Address(BaseModel):
    model_config = ConfigDict(strict=True, extra="forbid")

    street: str = Field(min_length=1, max_length=120)
    city: str = Field(min_length=1, max_length=60)
    country: CountryCode


class PackageLine(BaseModel):
    model_config = ConfigDict(strict=True, extra="forbid")

    sku: Sku
    quantity: int = Field(ge=1, le=10_000)
    # Decimal avoids binary float money errors. JSON has no decimal literal,
    # so this field opts OUT of strict mode (strict=False) and is hardened
    # by a validator instead -- blanket strictness breaks legitimate input.
    unit_weight_kg: Decimal = Field(strict=False, gt=Decimal("0"), le=Decimal("500"))

    @field_validator("quantity")
    @classmethod
    def quantity_not_suspiciously_round(cls, value: int) -> int:
        # Business heuristic: quantities over 5000 require explicit review.
        if value > 5000:
            raise ValueError("quantity above 5000 requires manual review")
        return value

    @field_validator("unit_weight_kg")
    @classmethod
    def weight_is_finite(cls, value: Decimal) -> Decimal:
        if not value.is_finite():
            raise ValueError("unit_weight_kg must be a finite number")
        return value.quantize(Decimal("0.001"))


class Shipment(BaseModel):
    """Strict at the boundary: no silent coercion, no unknown fields."""

    model_config = ConfigDict(strict=True, extra="forbid")

    shipment_id: Optional[str] = Field(default=None, pattern=r"^shp_\d{8}$")
    destination: Address
    lines: list[PackageLine] = Field(min_length=1)
    priority: str = Field(default="standard")

    @field_validator("priority")
    @classmethod
    def priority_is_known(cls, value: str) -> str:
        allowed = {"standard", "express", "freight"}
        if value not in allowed:
            raise ValueError(f"priority must be one of {sorted(allowed)}")
        return value

    @model_validator(mode="after")
    def express_shipments_stay_light(self) -> "Shipment":
        # Cross-field invariant: express parcels cannot exceed 30 kg total.
        if self.priority == "express":
            total = sum(
                line.unit_weight_kg * line.quantity for line in self.lines
            )
            if total > Decimal("30"):
                raise ValueError(
                    f"express shipments must weigh <= 30 kg, got {total} kg"
                )
        return self


def demo_coercion_pitfall() -> dict[str, str]:
    """Show why lax mode is dangerous at the API boundary."""
    outcome: dict[str, str] = {}

    class LaxPackage(BaseModel):  # default: lax coercion
        quantity: int

    coerced = LaxPackage(quantity="12")  # type: ignore[arg-type]
    outcome["lax_accepts_string"] = f"{coerced.quantity!r} (silently coerced)"

    try:
        PackageLine(sku="NRW-AR-1001", quantity="12", unit_weight_kg=Decimal("1.0"))  # type: ignore[arg-type]
    except Exception as exc:
        first = str(exc).splitlines()[0]
        outcome["strict_rejects_string"] = first
    return outcome


if __name__ == "__main__":
    for key, value in demo_coercion_pitfall().items():
        print(f"{key}: {value}")
\end{minted}

Running \texttt{python shipment\_models.py} prints:

\begin{minted}{text}
lax_accepts_string: 12 (silently coerced)
strict_rejects_string: 1 validation error for PackageLine
\end{minted}

The lax model accepted a string where an integer was declared, without a
murmur; the strict model refused. In lax mode the corruption would have
entered the system wearing a disguise; in strict mode it became a
\texttt{422} with pointer \texttt{/lines/0/quantity} that the client can fix.

\subsection{Redaction Middleware and Correlation IDs}

Listing~7.3 is the last line of defense. Two middlewares bracket the
application: \texttt{CorrelationIdMiddleware} stamps or mints the trace ID
(inner, so handlers can read \texttt{request.state.trace\_id}), and
\texttt{ProblemRedactionMiddleware} (outer) re-reads every
\texttt{application/problem+json} body and scrubs it. Middleware ordering in
Starlette is last-added-equals-outermost, so the redactor wraps the
correlator wraps the app. The redactor is deliberately pessimistic: if a
body smells like a stack trace, it replaces the detail entirely; the full
payload goes only to the server log, keyed by the trace ID, where operators
with access control can see it.

\begin{minted}{python}
"""redaction.py -- Correlation IDs and response redaction middleware.

Listing 7.3. Guarantees: problem+json bodies never leak PII patterns or
stack traces; full context is logged server-side under the trace_id.
"""
from __future__ import annotations

import json
import logging
import re
import uuid
from typing import Awaitable, Callable

from fastapi import Request, Response
from starlette.middleware.base import BaseHTTPMiddleware
from starlette.types import ASGIApp

logger = logging.getLogger("northwind.errors")

TRACE_HEADER = "X-Trace-Id"
PROBLEM_MEDIA_TYPE = "application/problem+json"

# Deterministic, conservative scrub rules. Extend via registry, not edits.
PII_PATTERNS: list[tuple[re.Pattern[str], str]] = [
    (re.compile(r"[\w.+-]+@[\w-]+\.[\w.]+"), "<email>"),
    (re.compile(r"\b\d{3}-\d{2}-\d{4}\b"), "<tax-id>"),
    (re.compile(r"\b(?:\d[ -]?){13,16}\b"), "<card-number>"),
    (re.compile(r"(?i)(api[_-]?key|token|password|secret)[\"']?\s*[:=]\s*[\"']?[^&\s\"']+"), r"\1=<redacted>"),
]

STACK_TRACE_MARKER = re.compile(r"(Traceback \(most recent call last\)|File \"[^\"]+\", line \d+)")


def scrub(text: str) -> str:
    """Remove PII-like patterns and stack-trace fragments from text."""
    for pattern, replacement in PII_PATTERNS:
        text = pattern.sub(replacement, text)
    if STACK_TRACE_MARKER.search(text):
        return "An internal error occurred. Contact support with the trace_id."
    return text


class CorrelationIdMiddleware(BaseHTTPMiddleware):
    """Accept or mint a trace ID and expose it on request.state + response."""

    async def dispatch(
        self,
        request: Request,
        call_next: Callable[[Request], Awaitable[Response]],
    ) -> Response:
        incoming = request.headers.get(TRACE_HEADER, "")
        # Accept only safe trace IDs: bounded length, no control chars.
        trace_id = incoming if re.fullmatch(r"[\w\-]{1,64}", incoming) else str(uuid.uuid4())
        request.state.trace_id = trace_id
        response = await call_next(request)
        response.headers[TRACE_HEADER] = trace_id
        return response


class ProblemRedactionMiddleware(BaseHTTPMiddleware):
    """Last line of defense: scrub every problem+json body before it leaves."""

    async def dispatch(
        self,
        request: Request,
        call_next: Callable[[Request], Awaitable[Response]],
    ) -> Response:
        response = await call_next(request)
        content_type = response.headers.get("content-type", "")
        if not content_type.startswith(PROBLEM_MEDIA_TYPE):
            return response

        body = b""
        async for chunk in response.body_iterator:  # type: ignore[attr-defined]
            body += chunk if isinstance(chunk, bytes) else chunk.encode()

        original = body.decode("utf-8", errors="replace")
        cleaned = scrub(original)
        if cleaned != original:
            trace_id = getattr(request.state, "trace_id", "unknown")
            # Server-side keeps the full payload for forensics; client gets scrubbed.
            logger.warning(
                "problem response redacted",
                extra={"trace_id": trace_id, "raw_body": original},
            )
        headers = dict(response.headers)
        headers["content-length"] = str(len(cleaned.encode()))
        return Response(
            content=cleaned,
            status_code=response.status_code,
            headers=headers,
            media_type=PROBLEM_MEDIA_TYPE,
        )
\end{minted}

\begin{warningbox}
\textbf{Regex redaction is a floor, not a ceiling.} Pattern lists catch the
shapes you remembered; they do not catch names, free-text addresses, or the
PII your schema did not anticipate. Treat regex scrubbing as the emergency
brake behind the primary defense---\emph{never putting PII into error
payloads in the first place}---and fuzz your scrubber with an adversarial
corpus (Capstone Project~3) before you trust it.
\end{warningbox}

\subsection{The Service, Wired}

Listing~7.4 assembles the pieces into a FastAPI application with an
in-memory store (offline, deterministic). It deliberately includes endpoints
for every failure mode we need to contract-test: missing resources, state
conflicts, a dependency failure, a throttle, and a ``leaky'' endpoint
simulating a developer who echoed PII into an error detail.

\begin{minted}{python}
"""app.py -- Northwind Robotics shipments service, wired for RFC 9457 errors.

Listing 7.4. In-memory store only: offline, deterministic, no network.
"""
from __future__ import annotations

from decimal import Decimal

from fastapi import FastAPI, Request

from problem_details import (
    ProblemException,
    RateLimitExceeded,
    ShipmentNotFound,
    ShipmentStateConflict,
    WarehouseUnavailable,
    register_problem_handlers,
)
from redaction import CorrelationIdMiddleware, ProblemRedactionMiddleware
from shipment_models import Shipment

# In-memory store keyed by shipment_id. Synthetic data only.
STORE: dict[str, Shipment] = {
    "shp_00000001": Shipment(
        shipment_id="shp_00000001",
        destination={
            "street": "1 Gearbox Lane",
            "city": "Springfield",
            "country": "US",
        },
        lines=[{"sku": "NRW-AR-1001", "quantity": 2, "unit_weight_kg": Decimal("4.5")}],
        priority="standard",
    )
}
SEQUENCE = {"next": 2}


def create_app() -> FastAPI:
    app = FastAPI(title="Northwind Robotics Shipments", version="1.0.0")
    # Middleware order matters: correlation first, redaction outermost-last.
    app.add_middleware(ProblemRedactionMiddleware)
    app.add_middleware(CorrelationIdMiddleware)
    register_problem_handlers(app)

    @app.post("/shipments", status_code=201)
    def create_shipment(shipment: Shipment, request: Request) -> Shipment:
        new_id = f"shp_{SEQUENCE['next']:08d}"
        SEQUENCE["next"] += 1
        stored = shipment.model_copy(update={"shipment_id": new_id})
        STORE[new_id] = stored
        return stored

    @app.get("/shipments/{shipment_id}")
    def get_shipment(shipment_id: str) -> Shipment:
        if shipment_id not in STORE:
            raise ShipmentNotFound(
                detail=f"No shipment with id {shipment_id!r} exists.",
                instance=f"/shipments/{shipment_id}",
            )
        return STORE[shipment_id]

    @app.post("/shipments/{shipment_id}/cancel")
    def cancel_shipment(shipment_id: str) -> Shipment:
        shipment = STORE.get(shipment_id)
        if shipment is None:
            raise ShipmentNotFound(detail="Cannot cancel a missing shipment.")
        if shipment.priority == "freight":
            # Domain rule: freight contracts lock at booking time.
            raise ShipmentStateConflict(
                detail="Freight shipments cannot be cancelled after booking.",
                instance=f"/shipments/{shipment_id}",
            )
        return shipment

    @app.get("/warehouse/ping")
    def warehouse_ping() -> dict[str, str]:
        # Simulated dependency failure: retriable, server-owned.
        raise WarehouseUnavailable(detail="Warehouse API timed out after 3 attempts.")

    @app.get("/shipments-throttled")
    def throttled() -> dict[str, str]:
        raise RateLimitExceeded(retry_after=30)

    @app.get("/debug/leaky")
    def leaky_failure() -> dict[str, str]:
        # A careless developer echoes request context into the error detail.
        raise ProblemException(
            detail="Failed to notify jane.doe@example.invalid; "
            "tax id 123-45-6789 on file."
        )

    return app


app = create_app()
\end{minted}

\subsection{Contract Tests: Every Error Path Is Schema-Valid}

Listing~7.5 is what turns the error contract from aspiration into CI gate.
Two ideas carry the weight. First, a \texttt{PROBLEM\_SCHEMA} (JSON Schema
draft 2020-12, \cite{jsonschemadocs}) with
\texttt{additionalProperties: false} that defines exactly what a problem
document from this service may contain---including the extension members.
Second, one assertion helper, \texttt{assert\_problem}, applied to
\emph{every} error test: correct status, correct media type, schema-valid
body, body status matching the header, and trace-ID round-trip. Adding a new
error path without a contract test becomes impossible to do quietly.

The tests run fully in-process via \texttt{TestClient}
\cite{fastapidocs}---no sockets, no network, deterministic on a fresh
machine. Note \texttt{raise\_server\_exceptions=False}: we want to verify
the \texttt{500} \emph{response}, not have the test harness re-raise.

\begin{minted}{python}
"""test_contract.py -- Contract tests: every error path is valid problem+json.

Listing 7.5. In-process ASGI tests via httpx TestClient. Fully offline.
"""
from __future__ import annotations

import jsonschema
import pytest
from fastapi.testclient import TestClient

from app import app

PROBLEM_SCHEMA: dict = {
    "$schema": "https://json-schema.org/draft/2020-12/schema",
    "type": "object",
    "required": ["type", "title", "status"],
    "properties": {
        "type": {"type": "string", "format": "uri-reference"},
        "title": {"type": "string", "minLength": 1},
        "status": {"type": "integer", "minimum": 400, "maximum": 599},
        "detail": {"type": "string"},
        "instance": {"type": "string"},
        "code": {"type": "string", "pattern": "^NRW-[A-Z-]+$"},
        "trace_id": {"type": "string", "minLength": 1},
        "errors": {
            "type": "array",
            "items": {
                "type": "object",
                "required": ["pointer", "code", "message"],
                "properties": {
                    "pointer": {"type": "string", "pattern": "^/"},
                    "code": {"type": "string"},
                    "message": {"type": "string"},
                },
            },
        },
    },
    "additionalProperties": False,
}

PII_FRAGMENTS = ("@", "123-45-6789", "Traceback", 'File "')


@pytest.fixture()
def client() -> TestClient:
    return TestClient(app, raise_server_exceptions=False)


def assert_problem(response, expected_status: int) -> dict:
    assert response.status_code == expected_status
    assert response.headers["content-type"].startswith("application/problem+json")
    body = response.json()
    jsonschema.validate(body, PROBLEM_SCHEMA)
    assert body["status"] == expected_status
    assert "X-Trace-Id" in response.headers
    assert body["trace_id"] == response.headers["X-Trace-Id"]
    return body


def test_404_is_problem_json(client: TestClient) -> None:
    body = assert_problem(client.get("/shipments/shp_99999999"), 404)
    assert body["code"] == "NRW-SHIPMENT-NOT-FOUND"
    assert body["type"].endswith("/shipment-not-found")


def test_422_maps_pydantic_errors_to_json_pointers(client: TestClient) -> None:
    payload = {
        "destination": {"street": "", "city": "Springfield", "country": "USA"},
        "lines": [{"sku": "bad-sku", "quantity": 0, "unit_weight_kg": "not-a-number"}],
        "priority": "hyperdrive",
    }
    body = assert_problem(client.post("/shipments", json=payload), 422)
    assert body["code"] == "NRW-VALIDATION-FAILED"
    pointers = {err["pointer"] for err in body["errors"]}
    # Each failure is located by an RFC 6901 JSON Pointer.
    assert "/destination/street" in pointers
    assert "/destination/country" in pointers
    assert "/lines/0/sku" in pointers
    assert "/lines/0/quantity" in pointers
    assert "/lines/0/unit_weight_kg" in pointers  # unparseable decimal
    assert "/priority" in pointers


def test_422_cross_field_invariant(client: TestClient) -> None:
    payload = {
        "destination": {"street": "1 Gearbox Lane", "city": "Springfield", "country": "US"},
        "lines": [{"sku": "NRW-AR-1001", "quantity": 10, "unit_weight_kg": 4.5}],
        "priority": "express",
    }
    body = assert_problem(client.post("/shipments", json=payload), 422)
    assert body["errors"][0]["pointer"] == "/"
    assert "express" in body["errors"][0]["message"]


def test_409_domain_conflict(client: TestClient) -> None:
    create = client.post(
        "/shipments",
        json={
            "destination": {"street": "9 Servo Rd", "city": "Springfield", "country": "US"},
            "lines": [{"sku": "NRW-FR-9000", "quantity": 1, "unit_weight_kg": 120}],
            "priority": "freight",
        },
    )
    assert create.status_code == 201
    shipment_id = create.json()["shipment_id"]
    body = assert_problem(client.post(f"/shipments/{shipment_id}/cancel"), 409)
    assert body["code"] == "NRW-SHIPMENT-STATE-CONFLICT"


def test_500_is_redacted_and_never_leaks_pii(client: TestClient) -> None:
    response = client.get("/debug/leaky")
    body = assert_problem(response, 500)
    for fragment in PII_FRAGMENTS:
        assert fragment not in response.text
    # The generic shape is preserved; the caller gets an actionable trace_id.
    assert body["code"] == "NRW-INTERNAL"


def test_429_carries_retry_after(client: TestClient) -> None:
    response = client.get("/shipments-throttled")
    assert_problem(response, 429)
    assert response.headers["Retry-After"] == "30"


def test_503_dependency_failure_is_retriable(client: TestClient) -> None:
    body = assert_problem(client.get("/warehouse/ping"), 503)
    assert body["code"] == "NRW-WAREHOUSE-UNAVAILABLE"


def test_client_supplied_trace_id_is_echoed(client: TestClient) -> None:
    response = client.get(
        "/shipments/shp_99999999", headers={"X-Trace-Id": "req-abc-123"}
    )
    assert response.headers["X-Trace-Id"] == "req-abc-123"
    assert response.json()["trace_id"] == "req-abc-123"


def test_malformed_trace_id_is_replaced(client: TestClient) -> None:
    response = client.get(
        "/shipments/shp_99999999", headers={"X-Trace-Id": "a" * 500}
    )
    trace_id = response.headers["X-Trace-Id"]
    assert trace_id != "a" * 500
    assert len(trace_id) <= 64


def test_happy_path_unaffected(client: TestClient) -> None:
    response = client.get("/shipments/shp_00000001")
    assert response.status_code == 200
    assert response.json()["shipment_id"] == "shp_00000001"
\end{minted}

Running the suite (all ten tests pass, in-process, in about a second):

\begin{minted}{text}
pytest -q
..........                                                               [100%]
10 passed in 1.05s
\end{minted}

\subsection{A Machine-Readable Error-Code Registry}

Listing~7.6 closes the loop between code, documentation, and operations. The
registry is a frozen dataclass table---one row per error code---that
\emph{validates itself at import time} (a duplicate code or missing
remediation crashes the deploy, not the customer) and renders the public
Markdown catalogue that the \texttt{type} URIs promise. Clients, docs, and
alert rules all derive from this single artifact; that is what
``machine-readable'' buys.

\begin{minted}{python}
"""error_codes.py -- Machine-readable error-code registry with doc generation.

Listing 7.6. One source of truth: clients, docs, and alerts all derive from
this registry. Run directly to print the generated Markdown catalogue.
"""
from __future__ import annotations

from dataclasses import dataclass
from enum import Enum


class Retriable(str, Enum):
    YES = "yes"
    NO = "no"
    CONDITIONAL = "conditional"


@dataclass(frozen=True)
class ErrorCode:
    code: str            # stable, machine-readable, never reused
    http_status: int
    type_slug: str       # suffix of the RFC 9457 type URI
    title: str           # short, stable, localizable key
    retriable: Retriable
    owner: str           # team accountable for this failure mode
    remediation: str     # one-line client-actionable guidance


REGISTRY: tuple[ErrorCode, ...] = (
    ErrorCode(
        code="NRW-VALIDATION-FAILED",
        http_status=422,
        type_slug="validation-failed",
        title="Validation Failed",
        retriable=Retriable.NO,
        owner="api-platform",
        remediation="Fix the fields named in errors[].pointer and resend.",
    ),
    ErrorCode(
        code="NRW-SHIPMENT-NOT-FOUND",
        http_status=404,
        type_slug="shipment-not-found",
        title="Shipment Not Found",
        retriable=Retriable.NO,
        owner="fulfillment",
        remediation="Verify the shipment_id; do not retry blindly.",
    ),
    ErrorCode(
        code="NRW-SHIPMENT-STATE-CONFLICT",
        http_status=409,
        type_slug="shipment-state-conflict",
        title="Shipment State Conflict",
        retriable=Retriable.CONDITIONAL,
        owner="fulfillment",
        remediation="Refresh the resource, re-evaluate preconditions, retry once.",
    ),
    ErrorCode(
        code="NRW-RATE-LIMIT-EXCEEDED",
        http_status=429,
        type_slug="rate-limit-exceeded",
        title="Rate Limit Exceeded",
        retriable=Retriable.YES,
        owner="api-platform",
        remediation="Back off for Retry-After seconds, then retry with jitter.",
    ),
    ErrorCode(
        code="NRW-WAREHOUSE-UNAVAILABLE",
        http_status=503,
        type_slug="warehouse-unavailable",
        title="Dependency Unavailable",
        retriable=Retriable.YES,
        owner="warehouse-integrations",
        remediation="Retry with exponential backoff; alert if it persists > 5 min.",
    ),
    ErrorCode(
        code="NRW-INTERNAL",
        http_status=500,
        type_slug="internal",
        title="Internal Server Error",
        retriable=Retriable.NO,
        owner="api-platform",
        remediation="Report the trace_id to support; do not retry mutations.",
    ),
)


def validate_registry(registry: tuple[ErrorCode, ...] = REGISTRY) -> None:
    """Fail fast at import time if the registry is internally inconsistent."""
    codes = [entry.code for entry in registry]
    assert len(codes) == len(set(codes)), "duplicate error codes"
    for entry in registry:
        assert entry.code.startswith("NRW-"), f"bad prefix: {entry.code}"
        assert 400 <= entry.http_status <= 599, f"bad status: {entry.code}"
        assert entry.remediation, f"missing remediation: {entry.code}"
        if entry.http_status < 500:
            # Client-facing codes must explain what the CLIENT can do.
            assert "support" not in entry.remediation.lower() or entry.http_status == 500


def render_markdown(registry: tuple[ErrorCode, ...] = REGISTRY) -> str:
    """Generate the public error catalogue page from the registry."""
    lines = [
        "# Northwind Robotics API -- Error Code Catalogue",
        "",
        "| Code | HTTP | Title | Retriable | Owner | Remediation |",
        "|---|---|---|---|---|---|",
    ]
    for entry in registry:
        lines.append(
            f"| `{entry.code}` | {entry.http_status} | {entry.title} "
            f"| {entry.retriable.value} | {entry.owner} | {entry.remediation} |"
        )
    return "\n".join(lines)


validate_registry()

if __name__ == "__main__":
    print(render_markdown())
\end{minted}

\section{Common Pitfalls \& Anti-Patterns}

Each entry below is a failure mode observed in real production estates, with
its cost and its cure.

\subsection*{200 OK with an Error Body}

\emph{Symptom:} \texttt{\{"success": false\}} rides a \texttt{200}.
\emph{Cost:} caches store failures, retry middleware sleeps through outages,
and availability dashboards show green during the fire. \emph{Cure:} the
status code reflects the outcome; the body explains it. No exceptions for
REST-style HTTP.

\subsection*{Leaking Stack Traces and PII}

\emph{Symptom:} a \texttt{500} body contains \texttt{Traceback (most recent
call last)}, a SQL fragment, or the caller's email address. \emph{Cost:}
reconnaissance data for attackers, regulatory exposure, PII smeared across
every log pipeline that captures bodies. \emph{Cure:} generic client-facing
\texttt{detail}, full context in server logs under the trace ID, redaction
middleware as the last gate (Listing~7.3).

\subsection*{Inconsistent Error Shapes per Endpoint}

\emph{Symptom:} \texttt{/orders} returns \texttt{\{"error": \ldots\}},
\texttt{/shipments} returns \texttt{\{"message": \ldots\}}, the gateway adds
an HTML page for good measure. \emph{Cost:} every consumer writes $N$ error
parsers and still misses the gateway's. \emph{Cure:} one pipeline, one media
type, schema-validated in CI for \emph{every} endpoint---including
framework-generated \texttt{404}s and \texttt{405}s.

\subsection*{Stringly-Typed Client Error Matching}

\emph{Symptom:} consumer code does
\texttt{if "insufficient" in response["message"]}. \emph{Cost:} your copy
edit is their breaking change; localization is a breaking change; fixing a
typo is a breaking change. \emph{Cure:} match on \texttt{type} and
\texttt{code} only; treat \texttt{title}/\texttt{detail} as display text.
Providers: never change \texttt{type}/\texttt{code} semantics in place.

\subsection*{500 for Validation Failures}

\emph{Symptom:} a bad payload raises \texttt{KeyError} deep in a handler and
surfaces as \texttt{500}. \emph{Cost:} clients give up (5xx reads as ``not
my fault, retry later''), your error budget burns for their typo, and your
on-call investigates non-incidents. \emph{Cure:} validate at the boundary
with a real schema; unreachable-by-validation code paths may still
\texttt{500}, and that is fine---that is what \texttt{500} is for.

\subsection*{Swallowing Exceptions into Generic 400}

\emph{Symptom:} \texttt{except Exception: return 400}. \emph{Cost:} the
mirror image of the previous sin: real outages masquerade as client errors,
clients stop retrying during your incident, and paging logic that keys on
5xx never fires. \emph{Cure:} catch specific exceptions, map each to its
true class, and let the genuinely unexpected fall through to the
\texttt{500} handler.

\subsection*{Logging Secrets from Failed Payloads}

\emph{Symptom:} a debug log line dumps the raw request body on validation
failure---including \texttt{password} and \texttt{api\_key} fields.
\emph{Cost:} credential material in log aggregators with broad read access
and long retention. \emph{Cure:} log the \emph{redacted} error summary, the
JSON Pointers of failing fields, and the trace ID---never the raw payload;
where payload capture is unavoidable, redact before writing (Chapter~5).

\section{Hands-On Production Lab \& Real-World Case Study}

\subsection{Lab: Retrofitting a Messy Service to Uniform RFC 9457}

You inherit \texttt{legacy\_quotes.py}, a FastAPI service for Northwind
Robotics' spare-parts quoting. It is a museum of error sins: one endpoint
returns \texttt{200} with \texttt{\{"ok": false\}}; another returns a bare
string; a third lets a \texttt{sqlite3.IntegrityError} escape as an HTML
\texttt{500}; validation errors use FastAPI's default (non-problem) shape;
and a ``helpful'' handler echoes the submitted customer email into its error
detail.

\textbf{Objective.} Without changing any success-path behavior, make every
error the service can emit conform to the RFC~9457 pipeline of
Section~\ref{sec:impl}.

\textbf{Steps.}

\begin{enumerate}
  \item \textbf{Inventory.} Exercise every endpoint's failure modes with
        \texttt{TestClient} and record the current (inconsistent) shapes in
        a table. This is your regression baseline.
  \item \textbf{Install the pipeline.} Drop in \texttt{problem\_details.py}
        and \texttt{redaction.py} unmodified; register handlers and
        middleware in the app factory.
  \item \textbf{Convert domain failures.} Replace ad-hoc returns with
        \texttt{ProblemException} subclasses; choose precise statuses
        (\texttt{409} for the duplicate-part-number case, \texttt{422} for
        semantic quote errors).
  \item \textbf{Fix the 200-with-error endpoint last}---it requires
        consumer coordination; document the change in the OpenAPI
        \texttt{responses} section \cite{openapi31}.
  \item \textbf{Write the contract tests first, per endpoint,} using the
        \texttt{PROBLEM\_SCHEMA} helper; run them red, then green.
  \item \textbf{Prove redaction.} Add the intentionally leaky handler back
        behind a test-only flag and assert no PII fragment survives.
\end{enumerate}

\textbf{Deliverable.} A diff touching no success-path logic, a contract test
per error path, and a one-page ``error contract'' section for the API's
README generated from your registry.

\subsection{War Story: The PII Leak That Hid in \texttt{detail}}

A fulfillment API (composite of several real incidents) validated customer
notification preferences against a downstream directory. When the directory
rejected an address, the developer---debugging an integration under deadline
pressure---returned the directory's error verbatim in the problem
\texttt{detail}: \emph{``mailbox jane.doe@example.invalid not provisioned for
account 123-45-6789.''} The response was a correctly-typed, correctly-shaped
\texttt{422}. Every control that looked at \emph{structure} passed.

The blast radius was invisible for weeks because the leak rode inside a
``valid'' error contract. Three systems made it worse: the CDN logged
response bodies for \texttt{422}s during a debugging flag that was never
turned off; the APM tool captured error bodies as ``context''; and a
third-party integrator helpfully surfaced \texttt{detail} in their
end-user-facing UI---so end users occasionally saw \emph{other} customers'
contact data when batch jobs interleaved failures.

Discovery came from an audit query, not an alert: a routine scan of APM
payloads for the regex \texttt{\textbackslash S+@\textbackslash S+} lit up
thousands of error bodies. Remediation proceeded in four moves, in this
order, and the order is the lesson: (1)~\emph{stop the bleeding at the
edge}---deploy redaction middleware the same day, before the code fix, because
the fix had other callers; (2)~\emph{fix the producer}---the directory
client now maps upstream errors to typed exceptions with generic details;
(3)~\emph{purge and rotate}---body capture disabled at the CDN, APM payload
retention shortened, affected logs expunged per policy; (4)~\emph{make it a
test, not a memory}---a contract test now asserts a PII-fragment denylist
against every error path, which is precisely
\texttt{test\_500\_is\_redacted\_and\_never\_leaks\_pii} in Listing~7.5.

The post-incident review concluded that the RFC~9457 adoption had been
\emph{structurally} complete and \emph{semantically} hollow: shape without
content policy. Every control the team added afterward---registry-reviewed
details, denylist contract tests, edge redaction---is about content. Shape is
necessary; it is nowhere near sufficient.

\section{Self-Assessment Exercises \& Deep-Dive Questions}

\begin{enumerate}
  \item A request with malformed JSON arrives at \texttt{POST /shipments}.
        Which layer of the taxonomy fires, what status results, and why does
        the answer \emph{not} involve Pydantic?
  \item Your teammate proposes returning \texttt{404} for unauthorized
        access to a resource. Construct the threat model where this is
        correct, and the one where it leaks information anyway.
  \item The registry in Listing~7.6 marks \texttt{409} as
        \emph{conditionally} retriable. Write the precise client algorithm
        (pseudo-code) for a safe \texttt{409} retry.
  \item Why does Listing~7.2 declare \texttt{unit\_weight\_kg} with
        \texttt{strict=False}? What breaks if you remove the opt-out, and
        what breaks if you also remove the companion validator?
  \item Design the JSON Pointer for a failure inside
        \texttt{lines[2].dimensions.width}. What RFC~6901 escaping must you
        apply if a key itself contains \texttt{/}?
  \item A client reports ``your API returns 500s.'' The dashboard shows
        5xx-dominated traffic that is actually \texttt{503} from your
        circuit breaker. What did the missing distinction cost you, and what
        metric label prevents recurrence?
  \item Argue both sides: should \texttt{title} be localized? What invariant
        must hold for localization to be safe?
  \item Your redaction middleware must not buffer 50~MB streaming problem
        bodies. Adapt Listing~7.3's design; what trade-off do you accept?
  \item Extend the contract suite so that a \emph{new} undocumented
        extension member fails CI. Which schema keyword already does this,
        and what workflow does that imply for adding members deliberately?
  \item Map the chapter's alert thresholds onto the error-budget policy of
        Chapter~18. Which single error code should \emph{never} page, and
        why?
\end{enumerate}

\section{Interview Q\&A Vault}

Thirty-four questions, junior through staff, with model answers calibrated to
what a strong candidate says out loud.

\subsection*{Junior}

\begin{enumerate}[label=\textbf{Q\arabic*.}, leftmargin=3.2em]
  \item \emph{What is the difference between \texttt{400} and \texttt{422}?}
        \textbf{A.} \texttt{400} means the request is malformed at the
        protocol/syntax level---unparseable JSON, bad framing. \texttt{422}
        means the request was syntactically valid but semantically
        unacceptable: the JSON parsed, the schema shape held, but the values
        violate domain rules. Mnemonic: \emph{400 = can't read you; 422 =
        read you, reject what you said.}
  \item \emph{What media type does RFC~9457 define, and why a distinct media
        type at all?} \textbf{A.} \texttt{application/problem+json}. The
        distinct type lets clients, gateways, and tooling identify error
        bodies without content sniffing, and lets the same endpoint serve a
        success schema and an error schema under one URL.
  \item \emph{Name the five standard members of a problem detail object.}
        \textbf{A.} \texttt{type} (URI identifying the problem class),
        \texttt{title} (stable short summary of the class), \texttt{status}
        (the HTTP status, echoed), \texttt{detail} (human explanation of
        this occurrence), \texttt{instance} (URI identifying this
        occurrence).
  \item \emph{Why is returning \texttt{200} with an error in the body
        harmful?} \textbf{A.} Every HTTP intermediary trusts the status
        code: caches may store the failure, retry logic sees success,
        monitoring reports availability during an outage. The status must
        reflect the outcome; the body elaborates.
  \item \emph{What is the difference between \texttt{401} and \texttt{403}?}
        \textbf{A.} \texttt{401} = unauthenticated (``prove who you are'')
        and must carry \texttt{WWW-Authenticate}; \texttt{403} =
        authenticated but not authorized (``I know you; no''). Re-login
        fixes the first, never the second.
  \item \emph{Should a client match on the \texttt{detail} string to decide
        handling?} \textbf{A.} Never. \texttt{detail} is human-facing,
        per-occurrence, and free to change (or be localized). Clients match
        on \texttt{type} and machine-readable extension members like
        \texttt{code}.
  \item \emph{What does \texttt{Retry-After} do and which statuses pair with
        it?} \textbf{A.} It tells the client how long to wait before
        retrying (delay-seconds or an HTTP date), and pairs with
        \texttt{429} and \texttt{503}. Without it, clients invent
        backoff---often synchronized, often harmful.
  \item \emph{In Pydantic v2, what is the difference between
        \texttt{field\_validator} and \texttt{model\_validator}?}
        \textbf{A.} A field validator sees one field and enforces local
        rules (pattern, range); a model validator (especially
        \texttt{mode="after"}) sees the whole validated model and enforces
        cross-field invariants like ``total weight under 30 kg for
        express.''
\end{enumerate}

\subsection*{Mid-Level}

\begin{enumerate}[label=\textbf{Q\arabic*.}, leftmargin=3.2em, start=9]
  \item \emph{What belongs in \texttt{detail} versus \texttt{title}?}
        \textbf{A.} \texttt{title} is per-\emph{type}: stable, short,
        localizable, safe to display as a heading; it must not vary between
        occurrences. \texttt{detail} is per-\emph{occurrence}: what happened
        this time, with the specifics a human needs---but never internals,
        secrets, or PII, and never anything a client would parse.
  \item \emph{How do you make errors machine-actionable?} \textbf{A.}
        Stable \texttt{type} URIs plus extension members: a short
        \texttt{code} for alerting and dashboards, a \texttt{retriable}
        classification (or per-code registry entry), structured field errors
        with JSON Pointers for \texttt{422}s, \texttt{trace\_id} for
        support, and remediation hints where the client can act.
  \item \emph{Why must \texttt{type} be a URI?} \textbf{A.} Global
        uniqueness without a central registry (mint under your domain),
        dereferenceability (the identifier doubles as the documentation
        link), and opacity (clients match the string, not fetched content).
        \texttt{about:blank} is the sanctioned generic fallback.
  \item \emph{When do you return \texttt{410} instead of \texttt{404}?}
        \textbf{A.} When you \emph{know} the resource existed and is
        permanently gone---deleted exports, sunset versions. \texttt{410} is
        cacheable and tells sync clients to purge. Use \texttt{404} when you
        do not know or do not want to reveal history; for sensitive
        resources, uniform \texttt{404} avoids an existence oracle.
  \item \emph{Give three distinct \texttt{409} conflict domains.}
        \textbf{A.} State-machine conflict (cancel a shipped shipment),
        uniqueness conflict (natural key taken), reference conflict (delete
        a warehouse holding inventory); a fourth is precondition/concurrency
        failure, though \texttt{412} is more precise for
        \texttt{If-Match}-style checks.
  \item \emph{Distinguish \texttt{500}, \texttt{502}, \texttt{503}, and
        \texttt{504} operationally.} \textbf{A.} \texttt{500}: my code
        faulted---page my team. \texttt{502}: upstream answered
        invalidly---page the integration owner. \texttt{503}: deliberate
        unavailability (shedding, maintenance, open breaker)---usually
        self-healing, pair with \texttt{Retry-After}. \texttt{504}: upstream
        never answered in time. The distinction drives paging, retry policy,
        and which dashboard lights up.
  \item \emph{Why is strict mode at the API boundary recommended, and what
        is the catch?} \textbf{A.} Strict mode rejects silent coercion
        (\texttt{"12"} $\rightarrow$ \texttt{12}), so client type bugs
        surface as fixable \texttt{422}s instead of corrupted data. The
        catch: JSON cannot express some Python types (\texttt{Decimal},
        \texttt{datetime} as literals), so those fields need per-field
        \texttt{strict=False} plus hardening validators.
  \item \emph{What is a JSON Pointer doing in an error response?}
        \textbf{A.} It locates the exact failing value in the submitted
        document (\texttt{/lines/0/quantity}, RFC~6901), so client tooling
        can highlight the offending field programmatically---unambiguous for
        machines, immune to message localization.
  \item \emph{Your service's default framework \texttt{404} returns HTML.
        Is that a problem?} \textbf{A.} Yes: it violates the uniform error
        contract---clients parsing problem+json get a surprise. Framework
        and gateway-generated errors (404, 405, 413) must pass through the
        same pipeline; contract tests should include at least one
        framework-generated error.
\end{enumerate}

\subsection*{Senior}

\begin{enumerate}[label=\textbf{Q\arabic*.}, leftmargin=3.2em, start=18]
  \item \emph{Design an error taxonomy for a payments API.} \textbf{A.}
        Three axes: layer (transport/protocol/application), ownership
        (client-fixable 4xx, server-owned 5xx, dependency-caused
        502/503/504 from the processor), and retriability---which for
        payments must compose with idempotency keys, because retrying a
        capture after a \texttt{504} without a key risks double-charging.
        Concrete classes: \texttt{422} semantic card data errors with
        per-field pointers; \texttt{402}-style declined payment as a
        distinct, documented type (not a generic 400); \texttt{409} for
        duplicate idempotency keys with different payloads; \texttt{503} +
        \texttt{Retry-After} for processor outages. Every class gets a
        registry row with owner, retriability, and remediation.
  \item \emph{How do you evolve an error contract without breaking
        clients?} \textbf{A.} Additively: new extension members are safe
        (clients must ignore unknowns); new \texttt{type} URIs for new
        classes are safe (clients must have a default branch). Never change
        the meaning of an existing \texttt{type}/\texttt{code}, never remove
        a documented member, and treat \texttt{title} copy edits as safe
        only because clients were told never to match on it.
  \item \emph{How do you keep PII out of error responses at scale?}
        \textbf{A.} Defense in depth: (1) type the pipeline so
        \texttt{detail} strings are built from curated templates, not
        exception messages; (2) never interpolate raw input values into
        messages; (3) redaction middleware as the last gate with a denylist
        of patterns and stack-trace markers; (4) contract tests asserting a
        PII-fragment denylist on every error path; (5) log-side redaction
        before writing failed payloads. Each layer assumes the previous one
        fails.
  \item \emph{How do you handle validation that requires a database---say,
        ``SKU must be stocked at destination''?} \textbf{A.} That is a
        domain invariant, not boundary validation: the payload is
        well-formed either way. Check it in the service layer and raise a
        typed domain exception (\texttt{409} or \texttt{422} with a specific
        \texttt{type} like \texttt{.../sku-not-stocked}). Do not smuggle
        database calls into Pydantic validators---models become untestable,
        and request latency becomes validation latency with no caching
        story.
  \item \emph{What should the correlation/trace ID story be across a
        multi-service request?} \textbf{A.} Mint or accept at the edge
        (validated: bounded length, character allowlist), propagate on every
        downstream call, embed in every problem body and log line, echo in
        response headers. Support flow: customer reports trace ID $\to$
        one query reconstructs the request tree. Never reflect unvalidated
        inbound IDs---header smuggling and log injection ride on that.
  \item \emph{A client asks you to add the offending \emph{value} to each
        \texttt{422} error for debugging convenience. Do you?} \textbf{A.}
        No as a default: the offending value is frequently a credential or
        PII pasted into the wrong field, and error bodies flow to logs,
        APMs, and client UIs. Offer the JSON Pointer and the rule; if a
        value echo is truly needed, gate it behind an authenticated
        debug-mode header that is disabled in production and redaction-aware.
  \item \emph{How do you test that your error contract holds, beyond unit
        tests?} \textbf{A.} Schema-validate every error path in contract
        tests with \texttt{additionalProperties: false} so undocumented
        members fail CI; property/fuzz test the redactor with an
        adversarial corpus; chaos-inject dependency failures and assert the
        mapped status family; and add a gateway-level test so intermediary
        errors match the contract too.
  \item \emph{When is \texttt{503} preferable to just letting requests time
        out?} \textbf{A.} Always, when you know you are shedding: a fast,
        explicit \texttt{503} + \texttt{Retry-After} preserves client
        goodwill, frees connections immediately, feeds accurate metrics, and
        lets well-behaved clients back off on schedule. A timeout costs the
        client the full deadline and teaches nothing.
  \item \emph{How does the error contract interact with caching?}
        \textbf{A.} Error statuses have distinct cacheability:
        \texttt{404}/\texttt{410} are cacheable by default (heuristic), and
        that is often desirable; \texttt{429}/\texttt{503} must not be
        cached beyond their validity---set \texttt{Cache-Control: no-store}
        on error responses that are occurrence-specific. A cached
        \texttt{500} replayed for an hour is an outage amplifier.
\end{enumerate}

\subsection*{Staff}

\begin{enumerate}[label=\textbf{Q\arabic*.}, leftmargin=3.2em, start=27]
  \item \emph{You are standardizing errors across forty services and six
        languages. What do you ship first?} \textbf{A.} The registry and the
        schema, not a library: a language-neutral problem schema, a governed
        error-code registry with owners and retriability, and a contract
        test kit each service can run in CI. Libraries per language follow,
        generated from or validated against those artifacts. Standardizing
        code before standardizing \emph{data} produces forty idiomatic
        dialects of inconsistency.
  \item \emph{How do you decide what is a new \texttt{type} versus a new
        \texttt{detail} of an existing type?} \textbf{A.} Ask whether a
        client should \emph{branch} on it. If handling differs---different
        remediation, different retriability---it is a new type. If only the
        human explanation differs, it is the same type with a different
        \texttt{detail}. Type proliferation is real: every type is a
        documented, versioned, supported promise.
  \item \emph{Design the alert policy on top of this chapter's taxonomy.}
        \textbf{A.} Metrics as \texttt{errors\_total\{endpoint, status,
        code\}}. Page on 5xx rate $>1\%$ over 5 min and on any single
        \texttt{code} spiking $>10\times$ baseline; ticket on sustained 4xx
        $>25\%$ (usually a broken consumer). Suppress \texttt{503}-shedding
        pages when the breaker is working as designed. Wire 5xx burn into
        the Chapter-18 error budget so deploys freeze on budget exhaustion.
  \item \emph{A regulator asks you to prove error responses never contain
        PII. What is your evidence chain?} \textbf{A.} Typed pipeline
        (details are template-built), middleware redaction (defense in
        depth), CI contract tests with PII denylists on every error path,
        a fuzzed redaction corpus with regression history, log-side
        redaction with access-controlled raw capture, and periodic audit
        scans of captured payloads. Process evidence (the pipeline) plus
        empirical evidence (the scans); neither alone is convincing.
  \item \emph{Should error responses be part of your OpenAPI document?
        What exactly do you publish?} \textbf{A.} Yes \cite{openapi31}: for
        each operation, the \texttt{4xx}/\texttt{5xx} responses with media
        type \texttt{application/problem+json}, the problem schema including
        your extension members, and \texttt{default} responses for the
        catch-alls. Per-operation, enumerate the \texttt{type} values that
        operation can emit---that enumeration is the actionable part of the
        contract.
  \item \emph{How do you keep gateway-generated errors (auth, throttling,
        body-size limits) inside the contract?} \textbf{A.} The gateway is
        part of the API as clients see it, so its errors must be mapped to
        problem+json at the edge: configure the gateway's error templates,
        mint gateway-scoped codes in the same registry
        (\texttt{NRW-GATEWAY-...}), and run contract tests through the
        gateway, not only against the app in-process.
  \item \emph{Your company acquires another firm whose API returns
        XML-only errors. Draft the migration.} \textbf{A.} Contract-first:
        publish the target problem+json schema and registry, stand up a
        translation shim at the edge that emits both (content negotiation on
        \texttt{Accept}), instrument which clients still negotiate XML,
        migrate them with sunset headers and a deadline, and keep
        \texttt{application/problem+xml} only until telemetry shows zero
        dependers. RFC~9457 retains the XML media type, so the shim is
        standards-legal throughout.
  \item \emph{What is the deepest mistake teams make with error design?}
        \textbf{A.} Treating errors as output formatting rather than as
        \emph{behavioral contract}: clients build retry loops, reconciliation
        jobs, and compliance evidence on your failures. The shape is the
        easy half; the hard half is semantics---ownership, retriability,
        stability guarantees, and the discipline to keep content policy
        (no PII, no internals) as tested code rather than wiki aspiration.
        Teams that get this right page less, integrate faster, and survive
        audits; teams that do not are debugging other people's string
        matching forever.
\end{enumerate}

\section{External Bibliography \& Further Reading}

\begin{itemize}
  \item \textbf{RFC 9457 --- Problem Details for HTTP APIs} \cite{rfc9457}.
        The primary source. Read \S3 (members), \S4 (extension members), and
        the security considerations; note its obsolescence of RFC~7807 and
        what changed (extension handling clarifications, modern HTTP
        alignment).
  \item \textbf{RFC 7807 (historical).} The 2016 original. Still cited in
        older libraries and articles; read it only to recognize legacy
        references, and follow 9457 for anything new.
  \item \textbf{RFC 9110 --- HTTP Semantics, \S15 (Status Codes)}
        \cite{rfc9110}. The authoritative definition of every status code
        discussed here, including cacheability defaults and the
        \texttt{WWW-Authenticate}/\texttt{Retry-After} header requirements.
  \item \textbf{RFC 6901 --- JSON Pointer.} The path syntax used for
        per-field error locations, including the \texttt{\textasciitilde 0}/
        \texttt{\textasciitilde 1} escaping rules.
  \item \textbf{Pydantic v2 documentation} \cite{pydanticdocs}. Validators,
        strict mode, \texttt{Annotated} constraints, and the error taxonomy
        (\texttt{loc}/\texttt{type}/\texttt{msg}) our mapper consumes.
  \item \textbf{FastAPI documentation} \cite{fastapidocs}. Exception
        handlers, middleware, \texttt{RequestValidationError}, and
        \texttt{TestClient} for in-process contract tests.
  \item \textbf{OpenAPI 3.1 specification} \cite{openapi31}. Documenting
        \texttt{application/problem+json} responses per operation and the
        \texttt{default} response.
  \item \textbf{JSON Schema documentation} \cite{jsonschemadocs}. Draft
        2020-12 validation semantics used by the contract tests.
  \item \textbf{Stripe API error documentation.} An exemplary industrial
        error contract: stable \texttt{type}, per-language handling guides,
        and \texttt{402}-class payment-specific codes worth studying as a
        design reference.
  \item \textbf{Google AIP-193 (Errors).} A mature enterprise standard for
        error payloads and codes; compare its decisions (canonical codes,
        \texttt{ErrorInfo} details) with RFC~9457's.
  \item \textbf{Microsoft REST API Guidelines --- error response section.}
        A second industrial reference point, useful for the
        \texttt{error.code}/\texttt{error.details} nesting pattern.
  \item \textbf{Nygard, \emph{Release It!}} \cite{nygard2018releaseit}.
        Failure modes, circuit breakers, and why \texttt{503}-as-shedding is
        a feature of healthy systems.
\end{itemize}

\section{Capstone Projects}

Five projects, progressive in difficulty. All are offline-first,
deterministic, and use synthetic Northwind Robotics data only.

\subsection{Project 1 [junior] --- Problem-Details Retrofit}

\textbf{Objective.} Take the intentionally messy \texttt{legacy\_quotes.py}
from this chapter's lab (or any small FastAPI service you wrote for
Chapter~6) and retrofit uniform RFC~9457 errors without touching
success-path behavior.

\textbf{Inputs/Datasets.} The lab service; this chapter's
\texttt{problem\_details.py}; a synthetic parts catalogue (10 fictional
SKUs).

\textbf{Steps.}
\begin{enumerate}
  \item Inventory current error shapes with \texttt{TestClient}.
  \item Register the problem handlers; convert ad-hoc errors to
        \texttt{ProblemException} subclasses.
  \item Write one contract test per error path using the shared
        \texttt{PROBLEM\_SCHEMA}.
  \item Update the OpenAPI \texttt{responses} documentation for two
        endpoints.
\end{enumerate}

\textbf{Acceptance Criteria.}
\begin{itemize}
  \item[$\square$] Every error response (including framework \texttt{404})
        is schema-valid \texttt{application/problem+json}.
  \item[$\square$] No success-path test changed; all still pass.
  \item[$\square$] \texttt{422}s carry JSON-Pointer-located
        \texttt{errors[]}.
  \item[$\square$] \texttt{pytest -q} green from a fresh \texttt{.venv}.
\end{itemize}

\textbf{Stretch Goals.} Add \texttt{Accept-Language}-driven localization of
\texttt{title} for two locales; add a client helper that branches on
\texttt{code}.

\subsection{Project 2 [mid] --- Error-Code Registry \& Documentation Generator}

\textbf{Objective.} Build a standalone registry package for a (fictional)
multi-service Northwind Robotics platform and generate both human docs and a
machine-readable artefact from it.

\textbf{Inputs/Datasets.} Listing~7.6 as a seed; a list of 15--20 invented
error codes across three services (shipments, warehouse, billing).

\textbf{Steps.}
\begin{enumerate}
  \item Model the registry with self-validation at import (duplicates,
        status ranges, remediation presence, retriability vocabulary).
  \item Generate a Markdown catalogue per service and a single
        \texttt{errors.json} for client SDKs.
  \item Write a CI-style check: any exception class whose \texttt{code} is
        absent from the registry fails the build (introspect the exception
        hierarchy).
  \item Emit a changelog diff between two registry versions.
\end{enumerate}

\textbf{Acceptance Criteria.}
\begin{itemize}
  \item[$\square$] Import-time validation catches an injected duplicate and
        an injected missing remediation (tests prove both).
  \item[$\square$] Generated Markdown renders a complete catalogue; JSON
        validates against a schema you publish.
  \item[$\square$] Orphan-code check fails a deliberately broken service.
  \item[$\square$] Registry diff output is deterministic (stable ordering).
\end{itemize}

\textbf{Stretch Goals.} Per-code deep links resolvable as \texttt{type}
URIs; deprecation workflow (mark, never delete, never reuse).

\subsection{Project 3 [mid] --- Redaction Library with Fuzz Corpus}

\textbf{Objective.} Harden Listing~7.3's scrubber into a standalone,
tested library and attack it yourself before an attacker does.

\textbf{Inputs/Datasets.} A hand-built adversarial corpus: 50+ synthetic
strings mixing fake emails (\texttt{*@example.invalid}), fake tax IDs, fake
card numbers, fake stack traces, near-misses (valid text that merely
\emph{looks} like PII), and Unicode edge cases.

\textbf{Steps.}
\begin{enumerate}
  \item Extract \texttt{scrub} into a package with a pluggable rule list.
  \item Build the corpus as data files; every corpus entry asserts
        redact/not-redact expectations.
  \item Add a deterministic property-style sweep: seeded generation of
        pattern-like mutations around each rule.
  \item Measure and bound false-positive rate on a benign-text corpus.
\end{enumerate}

\textbf{Acceptance Criteria.}
\begin{itemize}
  \item[$\square$] 100\% of malicious corpus entries redacted.
  \item[$\square$] Documented, bounded false-positive rate on benign
        corpus; regressions fail CI.
  \item[$\square$] All generation seeded; suite is deterministic and
        offline.
  \item[$\square$] Rule list extensible without editing library code.
\end{itemize}

\textbf{Stretch Goals.} Streaming redaction for chunked bodies; structured
(JSON-aware) scrubbing that redacts by field name as well as by pattern.

\subsection{Project 4 [senior] --- Error-Observability Dashboard Specification}

\textbf{Objective.} Specify---and prototype offline---the error
observability layer for the shipments API: metrics, alert rules, and a
dashboard definition, all derived from the error-code registry.

\textbf{Inputs/Datasets.} A synthetic, seeded log generator producing a week
of request outcomes with known injected incidents (a dependency outage, a
bad client deploy, a throttling event).

\textbf{Steps.}
\begin{enumerate}
  \item Define the metric taxonomy (\texttt{errors\_total\{endpoint, status,
        code\}}) and exemplar queries (error rate per code, per endpoint).
  \item Encode the chapter's alert thresholds as executable rules over the
        synthetic dataset.
  \item Prototype an offline aggregator that replays the log stream and
        prints alert firing timelines.
  \item Write the dashboard spec: panels, drill-down from code to trace
        IDs, error-budget burn panel referencing Chapter~18.
\end{enumerate}

\textbf{Acceptance Criteria.}
\begin{itemize}
  \item[$\square$] Replay deterministically detects all three injected
        incidents with zero false pages on the baseline day.
  \item[$\square$] Alert rules are data (declarative config), not code.
  \item[$\square$] Dashboard spec traces every panel to a registry code.
  \item[$\square$] Burn-rate panel matches hand-computed values on a
        100-request worked example.
\end{itemize}

\textbf{Stretch Goals.} Burn-rate windows (fast/slow) per Chapter~18;
anomaly baseline per code using the synthetic history.

\subsection{Project 5 [staff] --- Chaos Error-Injection Test Harness}

\textbf{Objective.} Build an in-process harness that injects failures
(timeouts, malformed upstream payloads, connection resets, slow leaks) at
the seams of a two-service synthetic system, and proves the \emph{entire
chain}---status mapping, problem shape, redaction, correlation, metrics
labels---survives chaos.

\textbf{Inputs/Datasets.} A two-service toy (gateway service calling a
warehouse stub); a seeded fault schedule; this chapter's full pipeline
installed in both services.

\textbf{Steps.}
\begin{enumerate}
  \item Implement fault injectors as dependency overrides: deterministic
        schedule from a seed, one fault per run or scripted mixes.
  \item Assert, per fault: correct status family (502/504/503, never 500),
        schema-valid problem body, trace ID continuity across both
        services, PII denylist holds.
  \item Produce a fault $\times$ contract matrix showing which faults each
        test catches; eliminate blind spots.
  \item Write the runbook entry template each fault class maps to.
\end{enumerate}

\textbf{Acceptance Criteria.}
\begin{itemize}
  \item[$\square$] Every injected fault yields a schema-valid problem
        response with the correct status family---no raw 500s, no HTML.
  \item[$\square$] Trace IDs join both services' logs for every scenario.
  \item[$\square$] Full suite deterministic under a fixed seed; runs
        offline in under a minute.
  \item[$\square$] Matrix reviewed: every fault class has at least one
        failing-before/passing-after test story.
\end{itemize}

\textbf{Stretch Goals.} Composed faults (timeout \emph{plus} redaction
bypass attempt); a ``contract score'' metric across the matrix trending
over registry versions.

\section{Python Dependencies Appendix}

Everything in this chapter installs from the following
\texttt{requirements.txt}:

\begin{minted}{text}
fastapi>=0.110
pydantic>=2.7
jsonschema>=4.21
httpx>=0.27
pytest>=8
\end{minted}

Install on Windows with PowerShell:

\begin{minted}{text}
python -m venv .venv
.\.venv\Scripts\Activate.ps1
pip install -r requirements.txt
\end{minted}

\subsection{fastapi ($\geq$ 0.110)}

\textbf{Description.} The ASGI web framework used for the service, its
exception-handler hooks, and \texttt{TestClient}. \textbf{Why included.}
Chapter-6 continuity; its \texttt{RequestValidationError} carries the
Pydantic error array our \texttt{422} mapper consumes, and its handler
registration is the seam where the single pipeline attaches.
\textbf{Alternatives.} Starlette alone (lower-level, same machinery); Flask
or Django REST framework (sync ecosystems with their own exception
patterns---the RFC~9457 design transfers, the code does not).
\textbf{Installation.} \texttt{pip install "fastapi>=0.110"}.

\subsection{pydantic ($\geq$ 2.7)}

\textbf{Description.} Data-validation library powering both the boundary
models (\texttt{Shipment}) and the \texttt{ProblemDetail} response model.
\textbf{Why included.} Strict-mode enforcement, \texttt{field\_validator} /
\texttt{model\_validator}, \texttt{Annotated} constraints, and a structured
error array with \texttt{loc} tuples that map cleanly to JSON Pointers.
\textbf{Alternatives.} \texttt{dataclasses} plus hand-rolled checks (verbose,
error-prone); \texttt{attrs}+\texttt{cattrs} (capable, less integrated with
FastAPI); Marshmallow (mature but schema-parallel to your models).
\textbf{Installation.} \texttt{pip install "pydantic>=2.7"}.

\subsection{jsonschema ($\geq$ 4.21)}

\textbf{Description.} JSON Schema validation for Python, used to
contract-test every error response against \texttt{PROBLEM\_SCHEMA}
\cite{jsonschemadocs}. \textbf{Why included.} Turns ``our errors have a
consistent shape'' from a hope into a CI gate;
\texttt{additionalProperties: false} makes undocumented extension members a
build failure. \textbf{Alternatives.} Hand-written assertions per field
(brittle); validating via OpenAPI tooling (heavier; OpenAPI~3.1 embeds JSON
Schema 2020-12 anyway \cite{openapi31}). \textbf{Installation.}
\texttt{pip install "jsonschema>=4.21"}.

\subsection{httpx ($\geq$ 0.27)}

\textbf{Description.} The HTTP client underneath FastAPI's
\texttt{TestClient} and the client used in this book's consumption chapters.
\textbf{Why included.} Enables fully in-process ASGI tests---no sockets, no
network, deterministic---while exercising the real request/response stack.
\textbf{Alternatives.} \texttt{requests} (no ASGI transport; Chapter~2
discusses the trade-offs); \texttt{aiohttp} (different async model).
\textbf{Installation.} \texttt{pip install "httpx>=0.27"}.

\subsection{pytest ($\geq$ 8)}

\textbf{Description.} Test runner and fixture machinery for the contract
suite. \textbf{Why included.} Fixture-scoped \texttt{TestClient}, expressive
assertion output, and the ecosystem standard; the suite is ten plain test
functions with one shared assertion helper. \textbf{Alternatives.}
\texttt{unittest} (stdlib, more ceremony); \texttt{nose2} (effectively
unmaintained). \textbf{Installation.} \texttt{pip install "pytest>=8"}.

\section{Chapter Summary \& Key Takeaways}

\begin{itemize}
  \item \textbf{Errors are contract.} Design, version, document, and test
        the failure surface with the same rigor as the success surface.
  \item \textbf{Classify before you code.} Transport versus protocol versus
        application; client-fixable (4xx) versus server-owned (5xx) versus
        dependency-caused; retriable versus not---with a written decision
        procedure, not vibes.
  \item \textbf{RFC~9457 is the shape.} \texttt{application/problem+json}
        with \texttt{type}/\texttt{title}/\texttt{status}/
        \texttt{detail}/\texttt{instance} plus documented extension members;
        \texttt{type} is the stable, URI-identified, documentation-anchored
        class; RFC~7807 is history.
  \item \textbf{Status precision pays.} \texttt{400} (syntax) versus
        \texttt{422} (semantics); \texttt{404} versus \texttt{410};
        \texttt{401} versus \texttt{403} versus \texttt{407};
        \texttt{429}+\texttt{Retry-After}; \texttt{500} versus
        \texttt{502}/\texttt{503}/\texttt{504}; and never
        \texttt{200}-with-error.
  \item \textbf{One pipeline.} Domain exceptions $\rightarrow$ problem
        mapper $\rightarrow$ redaction $\rightarrow$ response, with
        correlation IDs and full-fidelity server-side logging. No endpoint
        improvises.
  \item \textbf{Strict at the boundary, invariant-aware in the domain.}
        Pydantic strict mode with deliberate per-field exceptions; model
        validators for cross-field rules; service layer for stateful
        invariants.
  \item \textbf{Content policy is code.} No internals, no PII, no secrets in
        error bodies or logs---enforced by typed templates, denylist
        contract tests, and redaction middleware, not by convention.
  \item \textbf{Operate the taxonomy.} Metrics per endpoint, status, and
        code; page on 5xx burn; feed it into the Chapter-18 error budget.
\end{itemize}

\section{Quick-Reference Card}

\subsection*{Status-Code Selection}

\begin{table}[h]
  \centering
  \small
  \begin{tabularx}{\textwidth}{l l X}
    \toprule
    \textbf{Code} & \textbf{Retriable} & \textbf{Use when} \\
    \midrule
    400 & no & malformed syntax/framing; unreadable request \\
    401 & after auth & missing/failed credentials; send \texttt{WWW-Authenticate} \\
    403 & no & authenticated, not authorized (or use 404 to hide existence) \\
    404 & no & not found, or existence must not be revealed \\
    407 & after auth & proxy demands its own authentication \\
    409 & conditional & conflict with current resource state; refresh then retry \\
    410 & no & known-permanently-gone; clients should purge \\
    412 & conditional & precondition (\texttt{If-Match} etc.) failed \\
    415 & no & unsupported \texttt{Content-Type} \\
    422 & no & syntactically valid, semantically invalid; include field pointers \\
    429 & yes + \texttt{Retry-After} & client exceeds rate policy \\
    500 & only if idempotent & your code faulted; page your team \\
    502 & yes & upstream returned an invalid response \\
    503 & yes + \texttt{Retry-After} & deliberate shedding/maintenance/open breaker \\
    504 & yes & upstream timed out \\
    \bottomrule
  \end{tabularx}
  \caption{Error status selection cheat sheet.}
\end{table}

\subsection*{Problem-Details Members}

\begin{table}[h]
  \centering
  \small
  \begin{tabularx}{\textwidth}{l X}
    \toprule
    \textbf{Member} & \textbf{Rule of thumb} \\
    \midrule
    \texttt{type} & URI, stable forever, documents the class; clients match on it \\
    \texttt{title} & short, stable per type, localizable, display-only \\
    \texttt{status} & must equal the response status \\
    \texttt{detail} & per-occurrence, human-only, no PII/internals, never parsed \\
    \texttt{instance} & identifies this occurrence (path or URN) \\
    \texttt{code} (ext.) & short machine code from the registry; for alerts and SDKs \\
    \texttt{trace\_id} (ext.) & correlation ID; the support conversation starter \\
    \texttt{errors[]} (ext.) & per-field failures with RFC 6901 JSON Pointers \\
    \bottomrule
  \end{tabularx}
  \caption{Problem-details member cheat sheet.}
\end{table}

\subsection*{Redaction Checklist}

\begin{itemize}
  \item[$\square$] \texttt{detail} built from curated templates, not raw exception text.
  \item[$\square$] No stack traces, SQL, hostnames, or file paths in any client body.
  \item[$\square$] No PII (emails, tax IDs, card numbers) echoed from payloads.
  \item[$\square$] No secrets (tokens, keys, passwords) in bodies \emph{or} logs.
  \item[$\square$] Redaction middleware is the last gate on every problem+json response.
  \item[$\square$] Full context logged server-side under \texttt{trace\_id}, access-controlled.
  \item[$\square$] Contract tests assert a PII-fragment denylist on every error path.
  \item[$\square$] Redactor fuzzed against an adversarial corpus; suite deterministic.
\end{itemize}

\section{Standardized Evidence Addendum}
\subsection{Benchmark Snapshot Template}
\begin{table}[h]
  \centering
  \small
  \begin{tabularx}{\textwidth}{l c c c c}
    \toprule
    \textbf{Config} & \textbf{p50 latency} & \textbf{p99 latency} & \textbf{Error rate} & \textbf{Throughput} \\
    \midrule
    Baseline & -- & -- & -- & 1.00x \\
    Candidate A & -- & -- & -- & -- \\
    Candidate B & -- & -- & -- & -- \\
    \bottomrule
  \end{tabularx}
  \caption{Minimum benchmark table to keep per chapter.}
\end{table}

\subsection{Request Trace Template}
\begin{itemize}
  \item \textbf{Request:} one representative API call for this chapter's topic.
  \item \textbf{Before:} behavior (headers, payload, latency, failure mode) with the naive approach.
  \item \textbf{After:} behavior with the technique in this chapter.
  \item \textbf{Result:} what changed in correctness, resilience, latency, and operability.
\end{itemize}

\subsection{Cost and Latency Budget Snapshot}
\begin{table}[h]
  \centering
  \small
  \begin{tabularx}{\textwidth}{l c c}
    \toprule
    \textbf{Stage} & \textbf{Latency budget} & \textbf{Cost driver} \\
    \midrule
    TLS / connection setup & -- & handshake round trips, keep-alive hit rate \\
    Request validation & -- & schema complexity, CPU per request \\
    Business logic and I/O & -- & downstream fan-out, query cost \\
    Serialization and egress & -- & payload size, compression ratio \\
    \bottomrule
  \end{tabularx}
  \caption{Operational budget template for this chapter's technique.}
\end{table}

\section{Configuration Preset Template}
\begin{minted}{text}
# api_policy.yaml
policy_version: "v1.0.0"
component: "<chapter_topic>"
timeout_ms: 0
retry_budget: 0
fallback_strategy: "fail_fast"
rollback_trigger: "error_budget_burn_gt_threshold"
\end{minted}

\section{CI Quality Gate Specification}
\begin{itemize}
  \item contract tests green against the pinned OpenAPI/AsyncAPI/Protobuf spec,
  \item no breaking schema changes without a version bump,
  \item error responses conform to RFC 9457 problem-details shape,
  \item benchmark deltas within approved regression bounds (latency and throughput),
  \item policy version and rollback plan present in release metadata.
\end{itemize}

\section{Incident Runbook Template}
\begin{enumerate}
  \item Detect regression from SLO burn-rate alerts or consumer reports.
  \item Scope impacted endpoints, versions, and consumer segments.
  \item Contain impact (rollback, feature flag, rate-limit tightening, or traffic shift).
  \item Diagnose with traces, structured logs, and contract diffs.
  \item Recover with validated fix and verified rollout procedure.
  \item Prevent recurrence via new regression tests and CI gates.
\end{enumerate}

\section{Cross-Chapter Decision Card}
\begin{table}[h]
  \centering
  \small
  \begin{tabularx}{\textwidth}{l X X}
    \toprule
    \textbf{Dimension} & \textbf{Choose this approach when} & \textbf{Prefer an alternative when} \\
    \midrule
    Use case & -- & -- \\
    Latency budget & -- & -- \\
    Operational cost & -- & -- \\
    Team maturity & -- & -- \\
    \bottomrule
  \end{tabularx}
  \caption{Decision card to complete per chapter.}
\end{table}

\section{ROI Model and Build-vs-Buy}
\begin{itemize}
  \item \textbf{Build when:} the capability is differentiating, compliance forbids vendors, or scale makes vendor pricing irrational.
  \item \textbf{Buy when:} the capability is commodity (auth, gateways, observability), time-to-market dominates, or staffing is thin.
  \item \textbf{Hidden costs to model:} on-call load, upgrade cadence, expertise ramp, data egress fees.
\end{itemize}

\section{Disaster Recovery Notes (RPO/RTO)}
\begin{itemize}
  \item Define recovery point objective (RPO): maximum acceptable data loss window.
  \item Define recovery time objective (RTO): maximum acceptable restoration time.
  \item Test restores regularly; an untested backup is not a backup.
\end{itemize}

